% !TEX root = ../../main.tex
\setchapterimage{}
\setchapterintro{从环同态和理想的对应出发，讨论环的直和、分式环与多项式环，并研究整环及域上多项式的性质。}
\chapter{环}
\label{p2:ch3}

\section{环的同态定理}

环的同态同时保持加法与乘法。它的核是理想，因此可以先用加法群的同态定理构造商群之间的对应，再验证这些对应保持乘法。本节不要求环交换，也不要求有单位元；理想均指双边理想。

先说明交与和在子环、理想中分别具有哪些封闭性质。对非空子集，检验它是子环，可以验证任意两个元素之差及其乘积仍在其中；理想则还须验证左右乘以母环中任意元素后仍在其中。

\begin{proposition}
若$H,N$为环$R$的子环，则$H\cap N$为子环；

若$N$为$R$的理想，则$H\cap N$为$H$的理想。

更进一步，若$N,H$都是$R$的理想，那么$H\cap N$为$R$的理想。


\end{proposition}
\begin{proof}
若$x,y\in H\cap N$，则$x-y,xy$同时属于$H$和$N$，故$H\cap N$为$R$的子环。

若$N$为$R$的理想，则对$x\in H\cap N$、$h\in H$，有$hx,xh\in N$；又因$H$为子环，$hx,xh\in H$。因此
\[
hx,xh\in H\cap N,
\]
即$H\cap N$为$H$的理想。

若$H,N$均为$R$的理想，则对任意$r\in R$和$x\in H\cap N$，$rx,xr$同时属于$H$与$N$，所以
\[
rx,xr\in H\cap N.
\]
故$H\cap N$为$R$的理想。


\end{proof}

\begin{definition}
设$H,N$为环$R$的子环，定义它们的和为
\[
H+N=\{a+b:a\in H,\ b\in N\}.
\]
\end{definition}
\begin{proposition}
	$H+N $是$R$的加法子群。若$N$为$R$的理想，则$H+N$为$R$的子环；若$H,N$均为$R$的理想，则$H+N$也是理想。
\end{proposition}
\begin{proof}
	由于$0=0+0\in H+N$，此集合非空。若$a_i\in H,b_i\in N$，则
\[
(a_1+b_1)-(a_2+b_2)=(a_1-a_2)+(b_1-b_2)\in H+N.
\]
因此$H+N$是加法子群。若$N$为理想，则
	\begin{align*}
		(a_1+b_1)(a_2+b_2)
		&=a_1a_2+a_1b_2+b_1a_2+b_1b_2\in H+N.
	\end{align*}
	这里$a_1a_2\in H$，其余三项均在$N$中，所以乘积属于$H+N$，故它为子环。若$H$也为理想，则对$r\in R$有
	\[
	r(a+b)=ra+rb\in H+N,\qquad (a+b)r=ar+br\in H+N,
	\]
	所以$H+N$为理想。
\end{proof}
\begin{example}
设$F$为域，令$R=M_2(F)$，
\[
H=\left\{\begin{pmatrix}0&0\\a&0\end{pmatrix}:a\in F\right\},\qquad
N=\left\{\begin{pmatrix}0&b\\0&0\end{pmatrix}:b\in F\right\}.
\]
$H,N$中的任意两个矩阵相乘均为零，且各自对减法封闭，所以都是子环。但取
\[
\begin{pmatrix}0&0\\1&0\end{pmatrix}\in H,\qquad
\begin{pmatrix}0&1\\0&0\end{pmatrix}\in N,
\]
它们相乘得到
\[
\begin{pmatrix}0&0\\1&0\end{pmatrix}
\begin{pmatrix}0&1\\0&0\end{pmatrix}
=\begin{pmatrix}0&0\\0&1\end{pmatrix}\notin H+N.
\]
因为$H+N$中每个矩阵的对角线元素均为零，故$H+N$不是子环。这说明不能从两个子环直接推出其和仍为子环。
\end{example}

\begin{lemma}\label{p2:ring-image-preimage}
设$\sigma:R\to R'$为环同态，$N=\ker\sigma$，$H$为$R$的子环，$H'$为$R'$的子环。则
\begin{align*}
\sigma^{-1}(\sigma(H))&=H+N,\\
\sigma(\sigma^{-1}(H'))&=H'\cap\sigma(R).
\end{align*}
特别地，第一式等于$H$当且仅当$N\subseteq H$；若$\sigma$满射，则第二式等于$H'$。
\end{lemma}
\begin{proof}
\textbf{第一式：}先证$\sigma^{-1}(\sigma(H))\subseteq H+N$。任取$x\in\sigma^{-1}(\sigma(H))$，则$\sigma(x)\in\sigma(H)$，所以存在$h\in H$使$\sigma(x)=\sigma(h)$。于是
\[
\sigma(x-h)=\sigma(x)-\sigma(h)=0,
\]
即$x-h\in N$。因此
\[
x=h+(x-h)\in H+N.
\]

再证$H+N\subseteq\sigma^{-1}(\sigma(H))$。任取$x\in H+N$，写作$x=h+n$，其中$h\in H$、$n\in N$。由于$\sigma(n)=0$，
\[
\sigma(x)=\sigma(h)+\sigma(n)=\sigma(h)\in\sigma(H),
\]
故$x\in\sigma^{-1}(\sigma(H))$。合并两个包含关系，得到第一式。此式不要求$\sigma$满射。

因为$N$是理想，$H+N$为子环；任一同时包含$H$和$N$的子环都包含它们的和。因此$\sigma^{-1}(\sigma(H))$正是包含$H$和$\ker\sigma$的最小子环。又因$N\subseteq H+N$，
\[
\sigma^{-1}(\sigma(H))=H
\quad\Longleftrightarrow\quad H+N=H
\quad\Longleftrightarrow\quad N\subseteq H.
\]

\textbf{第二式：}任取$y\in\sigma(\sigma^{-1}(H'))$，则存在$x\in\sigma^{-1}(H')$使$y=\sigma(x)$。由反像的定义，$\sigma(x)\in H'$，故
\[
y\in H'\cap\sigma(R).
\]
反过来，若$y\in H'\cap\sigma(R)$，则存在$x\in R$使$y=\sigma(x)$；又因$y\in H'$，有$x\in\sigma^{-1}(H')$，所以
\[
y=\sigma(x)\in\sigma(\sigma^{-1}(H')).
\]
这就证明第二式。若$\sigma$满射，$\sigma(R)=R'$，于是
\[
\sigma(\sigma^{-1}(H'))=H'\cap R'=H'.
\]
\end{proof}

\begin{theorem}\label{p2:ring-correspondence}
\textbf{环的对应定理：}设$\sigma:R\to R'$为满环同态，$N=\ker\sigma$。映射$H\mapsto\sigma(H)$给出$R$中包含$N$的子环与$R'$的子环之间的一一对应；在此对应下，理想与理想对应。
\end{theorem}
\begin{proof}
\textbf{$\sigma(H)$是子环：}由于$0\in H$，$0\in\sigma(H)$。任取$\sigma(a),\sigma(b)\in\sigma(H)$，有
\[
\sigma(a)-\sigma(b)=\sigma(a-b)\in\sigma(H),\qquad
\sigma(a)\sigma(b)=\sigma(ab)\in\sigma(H).
\]
所以$\sigma(H)$为子环。

\textbf{$H$是理想，则$\sigma(H)$是理想：}若$H$为理想，任取$r'\in R'$及$a'\in\sigma(H)$，由满射取$r\in R,a\in H$使$\sigma(r)=r'$、$\sigma(a)=a'$，于是
\[
r'a'=\sigma(ra)\in\sigma(H),\qquad a'r'=\sigma(ar)\in\sigma(H).
\]
故$\sigma(H)$为理想。

\textbf{子环$H'\subseteq R'$的完全反像$\sigma^{-1}(H')$为子环：}若$x,y\in\sigma^{-1}(H')$，则$\sigma(x),\sigma(y)\in H'$，所以
\[
\sigma(x-y)=\sigma(x)-\sigma(y)\in H',\qquad
\sigma(xy)=\sigma(x)\sigma(y)\in H'.
\]
因此$x-y,xy\in\sigma^{-1}(H')$。又$0\in H'$，故$\ker\sigma\subseteq\sigma^{-1}(H')$，特别地该反像非空，是包含$N$的子环。

\textbf{理想$H'\subseteq R'$的完全反像$\sigma^{-1}(H')$为理想：}若$H'$为理想，任取$r\in R$、$a\in\sigma^{-1}(H')$，则
\[
\sigma(ra)=\sigma(r)\sigma(a)\in H',\qquad
\sigma(ar)=\sigma(a)\sigma(r)\in H',
\]
故$\sigma^{-1}(H')$为理想。

由引理\ref{p2:ring-image-preimage}，当$N\subseteq H$且$\sigma$满射时，
\[
\sigma^{-1}(\sigma(H))=H+N=H,\qquad
\sigma(\sigma^{-1}(H'))=H'.
\]
因此$H\mapsto\sigma(H)$与$H'\mapsto\sigma^{-1}(H')$互逆，得到所述一一对应。
\end{proof}

\begin{theorem}\label{p2:ring-third-iso}
设$\sigma:R\to R'$为环的满同态，$N=\ker\sigma$，$H$为$R$的任一包含$N$的理想。则$\sigma$诱导出环同构
\[
\overline\sigma:R/H\longrightarrow R'/\sigma(H),
\qquad
\overline\sigma(x+H)=\sigma(x)+\sigma(H).
\]
特别地，若取$R'=R/N$，$\sigma(x)=x+N$，则$\sigma(H)=H/N$，从而
\[
R/H\cong(R/N)/(H/N).
\]
同构映射为
\[
x+H\longmapsto(x+N)+(H/N).
\]
\end{theorem}
\begin{proof}
记$H'=\sigma(H)$。由环的对应定理，$H'$为$R'$的理想。由加法群的同构定理，$\sigma$诱导出加法群同构
\[
\overline\sigma:R/H\longrightarrow R'/H',
\qquad x+H\longmapsto\sigma(x)+H'.
\]

具体地，若$x+H=y+H$，则$x-y\in H$，故$\sigma(x)-\sigma(y)\in H'$，两者在$R'/H'$中的类相同，所以映射良好定义。它保持加法；由$\sigma$满射可知它满射。若$\overline\sigma(x+H)=H'$，则$\sigma(x)\in\sigma(H)$，于是
\[
x\in\sigma^{-1}(\sigma(H))=H+N=H.
\]
故核为零，确为加法群同构。这里用到了$N\subseteq H$，它正是单射性成立的条件。

再验证$\overline\sigma$保持乘法。任取$x,y\in R$，有
\begin{align*}
\overline\sigma\bigl((x+H)(y+H)\bigr)
&=\overline\sigma(xy+H)\\
&=\sigma(xy)+H'\\
&=\sigma(x)\sigma(y)+H'\\
&=(\sigma(x)+H')(\sigma(y)+H')\\
&=\overline\sigma(x+H)\,\overline\sigma(y+H).
\end{align*}
所以$\overline\sigma$是环同构。

取自然同态$\sigma:R\to R/N$，$\sigma(x)=x+N$，此时$\sigma(H)=H/N$。代入上述同构，即得特别情形及其同构映射。
\end{proof}

\begin{theorem}\label{p2:ring-second-iso}
若$H$为$R$的子环，$N$为$R$的理想，则
\[
H/(H\cap N)\cong(H+N)/N,\qquad
h+(H\cap N)\longmapsto h+N.
\]
\end{theorem}
\begin{proof}
将自然同态$\pi:R\to R/N$限制在$H$上。此限制的核恰为$\{h\in H:h\in N\}=H\cap N$。其像是$\pi(H)$，应用环同态基本定理得
\[
H/(H\cap N)\cong\pi(H).
\]
另一方面，$H+N$为子环且包含$N$，并有$\pi(H+N)=\pi(H)$。将$\pi$限制在$H+N$上，其核为$N$，于是
\[
(H+N)/N\cong\pi(H).
\]
合并两个同构得到结论；所列映射的逆映射由$h+n+N\mapsto h+(H\cap N)$给出。
\end{proof}

\begin{example}
设$m,n$为正整数，记$d=\gcd(m,n)$、$M=\operatorname{lcm}(m,n)$。在$\mathbb Z$中，
\[
(m)+(n)=(d),\qquad(m)\cap(n)=(M).
\]
第一式的理由是：和中的每个元素$am+bn$均被$d$整除；反过来，贝祖等式给出$d=um+vn$，所以每个$d$的倍数均在和中。第二式的理由是：同时属于两个理想，等价于同时被$m,n$整除，也就是被最小公倍数$M$整除。

在第二同构定理中取$H=(n)$、$N=(m)$，得到
\[
(n)/(M)\cong(d)/(m).
\]
左侧加法群有$M/n$个元素，右侧有$m/d$个元素：分别可取$0,n,\ldots,(M/n-1)n$及$0,d,\ldots,(m/d-1)d$为代表元。因此
\[
\frac Mn=\frac md,\qquad dM=mn.
\]
这也从理想与商环的关系说明最大公因数、最小公倍数之间的等式。
\end{example}


\section{环的直和}

环的直和与群的直和相平行：先作加法群的直和，再逐分量定义乘法。本节讨论的环默认有单位元，即为幺环；讨论任意环的理想生成时会另行说明。对幺环中的元素，零次幂记为$a^0=1$。

\begin{definition}
设$R_1,\ldots,R_r$为有限多个环。在它们的加法群直和上定义
\[
(a_1,\ldots,a_r)(b_1,\ldots,b_r)=(a_1b_1,\ldots,a_rb_r).
\]
所得环称为$R_1,\ldots,R_r$的\textbf{直和}，记为$R=R_1\oplus\cdots\oplus R_r$或$R=\bigoplus_{i=1}^rR_i$。
\end{definition}
\begin{proof}[补充验证]
加法已按分量定义，构成交换群。若$a=(a_i),b=(b_i),c=(c_i)$，则各个$R_i$中的结合律与分配律给出
\begin{align*}
(ab)c&=((a_ib_i)c_i)_i=(a_i(b_ic_i))_i=a(bc),\\
a(b+c)&=(a_i(b_i+c_i))_i=(a_ib_i+a_ic_i)_i=ab+ac,\\
(a+b)c&=((a_i+b_i)c_i)_i=(a_ic_i+b_ic_i)_i=ac+bc.
\end{align*}
零元为$(0_{R_1},\ldots,0_{R_r})$，单位元为$(1_{R_1},\ldots,1_{R_r})$。若每个$R_i$交换，则$ab=(a_ib_i)_i=(b_ia_i)_i=ba$，故直和也交换。
\end{proof}

\begin{proposition}
在$R=R_1\oplus\cdots\oplus R_r$中，令
\[
R_i'=\{(0,\ldots,0,a_i,0,\ldots,0):a_i\in R_i\}.
\]
则各$R_i'$都是$R$的理想，$R_i'\cong R_i$，并满足
\[
R=R_1'+\cdots+R_r',\qquad
R_i'\cap\Bigl(\sum_{j\ne i}R_j'\Bigr)=(0).
\]
每个元素都可唯一写为各$R_i'$中元素的和；$i\ne j$时$R_i'R_j'=(0)$。
\end{proposition}
\begin{proof}[补充证明]
第$i$坐标的映射$a_i\mapsto(0,\ldots,a_i,\ldots,0)$保持加法、乘法，且为到$R_i'$的双射，故给出同构。对$y=(b_1,\ldots,b_r)\in R$及$x=(0,\ldots,a_i,\ldots,0)\in R_i'$，
\[
yx=(0,\ldots,b_ia_i,\ldots,0),\qquad
xy=(0,\ldots,a_ib_i,\ldots,0),
\]
都属于$R_i'$，所以它是理想。

每个$(a_1,\ldots,a_r)$是各个仅含一个非零坐标的元素之和；第$i$个分量唯一决定第$i$项，所以表示唯一。若一个元素同时属于$R_i'$和其他坐标子环的和，它的所有坐标都为0，故交为$(0)$。不同坐标中的元素相乘，各分量至少有一个因子为0，故乘积为0。
\end{proof}

直和中的这些理想具有独立的加法分解和相互为零的乘法。下面反过来从环内部的理想恢复这种结构。

\begin{theorem}\label{p2:ring-internal-sum}
设$R_1,\ldots,R_r$为环$R$的理想，并满足
\[
R=R_1+\cdots+R_r,\qquad
R_i\cap(R_1+\cdots+\widehat{R_i}+\cdots+R_r)=(0).
\]
其中$\widehat{R_i}$表示在和中去掉$R_i$。则求和映射给出环同构
\[
R_1\oplus\cdots\oplus R_r\longrightarrow R,
\qquad(a_1,\ldots,a_r)\longmapsto a_1+\cdots+a_r.
\]
此时称$R$为这些子环的\textbf{内直和}。
\end{theorem}
\begin{proof}[补充证明]
由和等于$R$，每个元素都有这样的表示。若$\sum_i a_i=\sum_i b_i$，则对每个$i$，
\[
a_i-b_i=-\sum_{j\ne i}(a_j-b_j)
\in R_i\cap\sum_{j\ne i}R_j=(0),
\]
故$a_i=b_i$，表示唯一。所以求和映射是加法群同构。

若$i\ne j$，$a_i\in R_i,b_j\in R_j$，由理想的左右吸收性质，$a_ib_j$同时属于$R_i$和$R_j$。由于$R_i\cap R_j=(0)$，有$a_ib_j=0$。于是
\[
\Bigl(\sum_i a_i\Bigr)\Bigl(\sum_j b_j\Bigr)=\sum_i a_ib_i,
\]
正好与外直和的分量乘法一致。因此该加法群同构也保持乘法，是环同构。
\end{proof}

定理中的映射是求和映射。各$R_i$本身的单位元可以与$R$的单位元不同；这里的子环不要求包含母环同一个单位元。事实上，将$1$按直和写成$1=e_1+\cdots+e_r$，$e_i\in R_i$。对$a\in R_i$，不同分量的乘积为零，所以
\[
a=a1=ae_i,\qquad a=1a=e_i a.
\]
故$e_i$是$R_i$自身的单位元。这也解释了为什么这些理想仍可作为幺环参与直和。

为了给出与内直和判别相对应的商环形式，先定义理想的积与互素性。

\begin{definition}
设$H,N$为$R$的理想。全部有限和
\[
\sum_{i=1}^m a_ib_i,\qquad a_i\in H,\ b_i\in N,
\]
组成的集合称为$H,N$的\textbf{积}，记为$HN$，其中允许取零和。若$H+N=R$，则称$H,N$\textbf{互素}。
\end{definition}
\begin{proposition}
$HN$仍为$R$的理想，且对任意理想$K$，
\[
H(N+K)=HN+HK,\qquad(N+K)H=NH+KH.
\]
\end{proposition}
\begin{proof}[补充证明]
有限和相加仍为同类有限和，取负只须将每个$a_i$换成$-a_i$，所以$HN$是加法子群。对$r\in R$，
\[
r\sum_i a_ib_i=\sum_i(ra_i)b_i\in HN,
\qquad\Bigl(\sum_i a_ib_i\Bigr)r=\sum_i a_i(b_ir)\in HN.
\]
所以$HN$为双边理想。

对$h\in H,n\in N,k\in K$，$h(n+k)=hn+hk$，故$H(N+K)\subseteq HN+HK$。反过来，$N,K\subseteq N+K$说明$HN,HK\subseteq H(N+K)$，两者之和也包含于其中。第二式同样由$(n+k)h=nh+kh$及两个反向包含得到。
\end{proof}

\begin{lemma}\label{p2:ring-comaximal-product}
设$H,N,K$为幺环$R$的理想。
\begin{enumerate}
\item 若$R$交换且$H+N=R$，则$HN=H\cap N$。
\item 若$H+N=R$且$K+N=R$，则$HK+N=R$。
\end{enumerate}
\end{lemma}
\begin{proof}
\textbf{第一条：}任一乘积$hn$同时属于$H,N$，所以总有$HN\subseteq H\cap N$。由互素性，取$a\in H,b\in N$使$a+b=1$。若$c\in H\cap N$，则
\[
c=c(a+b)=ca+cb=ac+cb.
\]
其中$ac\in HN$，因为$a\in H,c\in N$；$cb\in HN$，因为$c\in H,b\in N$。所以$c\in HN$，得到反向包含。这里把$ca$写成$ac$使用了交换性。

\textbf{第二条：}取$a\in H,b\in K,c,d\in N$使$a+c=1$、$b+d=1$。于是
\[
1=(a+c)(b+d)=ab+(ad+cb+cd)\in HK+N.
\]
一个含$1$的理想包含每个$x=x1$，所以$HK+N=R$。这一步不要求$R$交换。
\end{proof}

\begin{theorem}\label{p2:ring-crt}
设幺环$R$的理想$N_1,\ldots,N_r$两两互素，$\sigma_i:R\to R/N_i$为自然同态。则
\[
\sigma:R\longrightarrow\bigoplus_{i=1}^rR/N_i,
\qquad x\longmapsto(\sigma_1(x),\ldots,\sigma_r(x))
\]
为满环同态，且$\ker\sigma=\bigcap_iN_i$。因此
\[
R/(N_1\cap\cdots\cap N_r)\cong R/N_1\oplus\cdots\oplus R/N_r.
\]
\end{theorem}
\begin{proof}
$r=1$时即自然同态的结论，以下设$r\geq2$。每个分量保持运算，所以$\sigma$为环同态。关键是证明它满射。

记$M_i=N_1\cdots N_{i-1}N_{i+1}\cdots N_r$，各理想按原次序相乘。理想积的分配律给出
\begin{align*}
M_1+M_2&=(N_2+N_1)N_3\cdots N_r=N_3\cdots N_r,\\
M_1+M_2+M_3&=(N_3+N_1N_2)N_4\cdots N_r=N_4\cdots N_r.
\end{align*}
第二个等式用到前一引理：$N_1,N_2$都与$N_3$互素，所以$N_1N_2$也与$N_3$互素。继续归纳，$N_1\cdots N_{k-1}$与$N_k$互素，故每加入$M_k$，就得到
\[
M_1+\cdots+M_k=N_{k+1}\cdots N_r.
\]
最后$k=r$时右侧按空乘积理解为理想$R$，所以$\sum_iM_i=R$。

因此可取$e_i\in M_i$满足$e_1+\cdots+e_r=1$。当$j\ne i$时，$M_j\subseteq N_i$，故
\[
\sigma_i(e_j)=0\quad(j\ne i),\qquad\sigma_i(e_i)=1+N_i.
\]
任给直和中的元素$(x_1+N_1,\ldots,x_r+N_r)$，令$x=\sum_j e_jx_j$，则
\begin{align*}
\sigma_i(x)&=\sum_j\sigma_i(e_j)\sigma_i(x_j)\\
&=\sigma_i(e_i)\sigma_i(x_i)=\sigma_i(x_i).
\end{align*}
所以$\sigma(x)$就是任给的元素，$\sigma$满射。

最后，$\sigma(x)=0$当且仅当每个$\sigma_i(x)=0$，即$x\in N_i$对所有$i$成立。因此核为$\bigcap_iN_i$，再用环同态基本定理得到商环同构。
\end{proof}

\begin{definition}
设$N\lhd R$。若$a-b\in N$，则称$a,b$\textbf{模$N$同余}，记为$a\equiv b\pmod N$。它等价于$a+N=b+N$。
\end{definition}

同余类就是商环的元素，所以同余保持加法和乘法。例如$ax\equiv b\pmod N$等价于商环中的等式$(a+N)(x+N)=b+N$。

\begin{corollary}\label{p2:ring-crt-congruence}
\textbf{中国剩余定理：}若$N_1,\ldots,N_r$两两互素，则对任给$b_1,\ldots,b_r\in R$，同余方程组
\[
x\equiv b_i\pmod{N_i}\qquad(1\leq i\leq r)
\]
有解，且任意两解模$N_1\cap\cdots\cap N_r$同余。
若该交为$(0)$，则$R\cong\bigoplus_iR/N_i$。
\end{corollary}
\begin{proof}[补充证明]
存在性是上一定理中$\sigma$的满射性。若$x,y$均为解，则$x-y\in N_i$对所有$i$成立，所以$x-y\in\bigcap_iN_i$；反之，在一个解上加上交中任意元素仍是解。因此解模该交唯一。交为零时，$\sigma$也单射，得到最后的同构。
\end{proof}

两种直和描述可以互相转化。若交为零且两两互素，令$R_i=\bigcap_{j\ne i}N_j$，在上述同构下，$R_i$恰对应只有第$i$个坐标可能非零的理想，故$R$是这些$R_i$的内直和。反过来，若$R$已是$R_1,\ldots,R_r$的内直和，令$N_i=\sum_{j\ne i}R_j$，则$N_i$两两互素、交为零，且$R/N_i\cong R_i$。这是同一分解的两种表达方式。

\begin{example}
在$\Z$中，设$n=\prod_{i=1}^r p_i^{e_i}$，各$p_i$为不同素数。由贝祖等式，理想$(p_i^{e_i})$两两互素。反复应用$HN=H\cap N$，得到它们的交为$(n)$，所以
\[
\Z/(n)\cong\bigoplus_{i=1}^r\Z/(p_i^{e_i}).
\]
\end{example}

本节最后讨论任意环中的理想生成，此处不要求环有单位元。

\begin{definition}
对非空子集$S\subseteq R$，所有包含$S$的理想之交称为\textbf{由$S$生成的理想}，记为$(S)$。它是包含$S$的最小理想。若$S=\{a_1,\ldots,a_r\}$，简记为$(a_1,\ldots,a_r)$，称为\textbf{有限生成理想}；由一个元素生成的理想$(a)$称为\textbf{主理想}。
\end{definition}

\begin{proposition}
任意环中，主理想$(a)$由全部形如
\[
na+xa+ay+\sum_{i=1}^m x_iay_i,
\qquad n\in\Z,\quad x,y,x_i,y_i\in R
\]
的元素组成。若$R$交换，可化为$na+xa$；若又有单位元，则$(a)=Ra$。此外，
\[
(a_1,\ldots,a_r)=(a_1)+\cdots+(a_r).
\]
\end{proposition}
\begin{proof}[补充证明]
令$J$为所列元素组成的集合。它对加法、取负封闭，且取$n=1$、其余项为零知$a\in J$。左乘$r\in R$时，四类项分别成为$(nr)a$、$(rx)a$、$ray$及$(rx_i)ay_i$；右乘时分别成为$a(nr)$、$xar$、$a(yr)$及$x_i a(y_ir)$，仍为所列项的有限和。因此$J$是含$a$的理想。

任一包含$a$的理想必须包含它的整数倍、左右乘积及这些项的有限和，故必包含$J$。所以$J=(a)$。交换时双侧项$xay$也可写为$(xy)a$，各项合并为$na+xa$；有单位元时$na=(n1)a$，再合并得到$Ra$。

最后，右侧的理想和含全部$a_i$，故包含$(a_1,\ldots,a_r)$；任何含全部$a_i$的理想又必须包含各$(a_i)$，所以包含它们的和。两边相等。
\end{proof}


%\section{环的反同构}
%
%环同构保持加法与乘法。矩阵转置提供另一种自然的对应：它保持加法，却把乘法的次序颠倒。对数域$F$上的矩阵环$M_n(F)$，
%\[
%(A+B)^T=A^T+B^T,\qquad(AB)^T=B^TA^T.
%\]
%这引出反同构的概念。
%
%\begin{definition}
%若双射$\sigma:R\to R'$满足
%\[
%\sigma(a+b)=\sigma(a)+\sigma(b),\qquad
%\sigma(ab)=\sigma(b)\sigma(a),
%\]
%则称$\sigma$为$R$到$R'$的\textbf{反同构}，并称$R,R'$\textbf{反同构}。若$R'=R$，则称为\textbf{反自同构}。不要求双射时，满足这两式的映射称为\textbf{反同态}。
%\end{definition}
%
%\begin{proposition}
%若$R,R'$反同构，且$R$交换，则$R'$也交换，这时反同构同时也是同构。
%\end{proposition}
%\begin{proof}[补充证明]
%设$u=\sigma(a),v=\sigma(b)$为$R'$中任意元素。由$ab=ba$，
%\[
%uv=\sigma(ba)=\sigma(ab)=vu.
%\]
%故$R'$交换，且$\sigma(ab)=\sigma(b)\sigma(a)=\sigma(a)\sigma(b)$，也保持通常次序的乘法。
%\end{proof}
%
%\begin{example}
%\textbf{矩阵转置：}对交换域$F$上的矩阵，转置两次得到原矩阵，所以它是双射。又
%\[
%((AB)^T)_{ij}=(AB)_{ji}=\sum_k a_{jk}b_{ki}
%=\sum_k b_{ki}a_{jk}=(B^TA^T)_{ij}.
%\]
%因此转置为$M_n(F)$的反自同构。上述计算使用系数交换性；在非交换系数环上，普通转置不一定颠倒乘法。
%\end{example}
%
%\begin{example}
%\textbf{四元数共轭：}四元数体$H$中的元素写为
%\[
%\alpha=a_0+a_1I+a_2J+a_3K,\qquad a_i\in\R,
%\]
%其中$I^2=J^2=K^2=-1$、$IJ=-JI=K$、$JK=-KJ=I$、$KI=-IK=J$。定义
%\[
%\overline\alpha=a_0-a_1I-a_2J-a_3K.
%\]
%则$\alpha\mapsto\overline\alpha$是$H$的反自同构，且共轭两次为恒等映射。
%\end{example}
%\begin{proof}
%共轭保持加法，且$\overline{\overline\alpha}=\alpha$，所以它双射。为验证$\overline{\alpha\beta}=\overline\beta\,\overline\alpha$，由实数线性及分配律，只须在基$1,I,J,K$的乘积上验证。
%
%若其中一项为1，等式由定义成立。若两项为同一个虚基元$U$，则
%\[
%\overline{U^2}=-1=(-U)(-U)=\overline U\,\overline U.
%\]
%若为不同虚基元$U,V$，则$UV$为另一个虚基元或其负数，且$VU=-UV$，所以
%\[
%\overline{UV}=-UV=VU=(-V)(-U)=\overline V\,\overline U.
%\]
%这涵盖所有基的乘积，再对任意四元数展开即可得到所需等式。
%\end{proof}
%
%\begin{example}
%\textbf{四元数矩阵的共轭转置：}对$A=(\alpha_{ij})\in M_n(H)$，令$\overline A=(\overline{\alpha_{ij}})$，再取转置$\overline A^{\,T}$。则$A\mapsto\overline A^{\,T}$是反自同构，连续作用两次得到$A$。
%\end{example}
%\begin{proof}[补充证明]
%记$A^*=\overline A^{\,T}$。它保持加法且$(A^*)^*=A$。对$A=(\alpha_{ij}),B=(\beta_{ij})$，
%\begin{align*}
%((AB)^*)_{ij}
%&=\overline{\sum_k\alpha_{jk}\beta_{ki}}\\
%&=\sum_k\overline{\beta_{ki}}\,\overline{\alpha_{jk}}\\
%&=(B^*A^*)_{ij}.
%\end{align*}
%所以$(AB)^*=B^*A^*$。共轭与转置结合起来，才在非交换的四元数矩阵环上得到反自同构。
%\end{proof}
%
%\begin{proposition}
%环$R$的全部自同构与反自同构在复合下组成一个群$G$。自同构群$\Aut(R)$是它的子群，且$[G:\Aut(R)]\leq2$。
%\end{proposition}
%\begin{proof}[补充证明]
%两个保持乘法次序的映射复合仍保持次序；一个保持次序、一个颠倒次序时，复合颠倒次序；两个都颠倒次序时，复合恢复原次序。例如反自同构$\sigma,\tau$满足
%\[
%(\sigma\tau)(ab)=\sigma(\tau(b)\tau(a))
%=(\sigma\tau)(a)(\sigma\tau)(b).
%\]
%逆映射具有与原映射相同的保序或反序性质，恒等映射也在其中，因此构成群。
%
%若不存在反自同构，或$R$交换使两种映射重合，则$G=\Aut(R)$。否则固定一个反自同构$\tau$，任一反自同构$\sigma$都使$\sigma\tau^{-1}$为自同构，所以全部反自同构组成一个陪集$\Aut(R)\tau$。故至多有两个陪集。
%\end{proof}
%
%即使一个非交换环没有反自同构，也总能构造另一个与它反同构的环。取与$R$一一对应的副本$R'=\{x':x\in R\}$，规定
%\[
%x'+y'=(x+y)',\qquad x'y'=(yx)'.
%\]
%加法群结构不变；乘法满足结合律，因为$(x'y')z'=(zyx)'=x'(y'z')$，分配律也由$R$的分配律得到。于是$R'$为环，而$x\mapsto x'$为反同构。这个构造并不表示$R$一定存在到自身的反自同构。
%
%
%\section{素理想和极大理想}
%
%从本节开始，除特别说明外，讨论的环均为交换幺环。域与整环均要求$1\ne0$。理想$N$给出商环$R/N$，下面研究$N$具有什么性质时，该商环成为域或整环。
%
%\begin{definition}
%若理想$P\ne R$满足
%\[
%ab\in P\quad\Longrightarrow\quad a\in P\text{ 或 }b\in P,
%\]
%则称$P$为\textbf{素理想}。若理想$M\ne R$，且不存在理想$A$使$M\subsetneq A\subsetneq R$，则称$M$为\textbf{极大理想}。
%\end{definition}
%
%\begin{example}
%$\Z$中的素理想恰为$(0)$及各素数$p$生成的$(p)$；极大理想恰为这些$(p)$。
%\end{example}
%\begin{proof}[补充证明]
%若$ab\in(p)$，则$p\mid ab$，由$p$为素数得$p\mid a$或$p\mid b$，故$(p)$为素理想。若$(p)\subseteq(n)\subseteq\Z$，则$n\mid p$，取$n\geq0$后只能$n=1$或$p$，所以$(p)$极大。
%
%$(0)$为素理想，因为$\Z$无零因子；它不极大，例如$(0)\subsetneq(2)\subsetneq\Z$。其余非零真理想为$(n)$、$n>1$。若$n$合数，写$n=ab$、$1<a,b<n$，则$ab\in(n)$而$a,b\notin(n)$，故不是素理想；取$n$的素因子$p$又有$(n)\subsetneq(p)\subsetneq\Z$，所以也不是极大理想。
%\end{proof}
%
%\begin{theorem}\label{p2:prime-maximal}
%设$R$为交换幺环。
%\begin{enumerate}
%\item 满环同态$\sigma:R\to R'$诱导包含$\ker\sigma$的极大理想与$R'$的极大理想之间的一一对应。
%\item $R$是域，当且仅当$(0)$为极大理想。
%\item 理想$M\ne R$极大，当且仅当$R/M$是域。
%\end{enumerate}
%\end{theorem}
%\begin{proof}
%\textbf{第一条：}环的理想对应保持并反映包含关系，且把$R$映到$R'$。若包含核的真理想$M$与$R$之间存在中间理想，则其像也位于$\sigma(M)$与$R'$之间；反过来，目标中的中间理想的完全反像也给出源中的中间理想。所以极大性在对应下保持。
%
%\textbf{第二条：}若$R$为域，任一非零理想$H$含非零元$a$，于是也含$a^{-1}a=1$，从而$H=R$。所以$(0)$极大。
%反过来，若$(0)$极大，对任意$a\ne0$，$Ra$是含$a$的非零理想，故$Ra=R$。于是存在$c\in R$使$ca=1$。交换性给出$ac=1$，所以每个非零元可逆，$R$为域。
%
%\textbf{第三条：}把第一条用于自然满同态$R\to R/M$，$M$对应零理想。因此$M$极大当且仅当$R/M$的零理想极大，再由第二条即得结论。
%\end{proof}
%
%\begin{theorem}\label{p2:prime-domain}
%理想$P\ne R$为素理想，当且仅当$R/P$为整环。特别地，$R$为整环当且仅当$(0)$为素理想。满环同态$\sigma:R\to R'$还诱导包含$\ker\sigma$的素理想与$R'$的素理想之间的一一对应。
%\end{theorem}
%\begin{proof}
%在$R/P$中，用$\overline a=a+P$表示剩余类。若$P$为素理想，则
%\[
%\overline a\,\overline b=0
%\Longrightarrow ab\in P
%\Longrightarrow a\in P\text{ 或 }b\in P
%\Longrightarrow\overline a=0\text{ 或 }\overline b=0.
%\]
%又$P\ne R$保证$R/P$非零，所以它是整环。
%反过来，若商环为整环，则$ab\in P$使$\overline a\,\overline b=0$，故至少一个类为零，即至少一个因子属于$P$。取$P=(0)$便得特别情形。
%
%最后，若$\ker\sigma\subseteq P$，由环同态定理，$R/P\cong R'/\sigma(P)$。两商环同时为整环或同时不是整环，结合理想对应，得到素理想的对应。
%\end{proof}
%
%域都是整环，所以每个极大理想均为素理想。反过来不成立，例如$\Z$的零理想。下面证明，在适当避开指定元素的条件下，素理想一定存在。
%
%\begin{definition}
%若$a\in R$满足$a^m=0$对某个正整数$m$成立，则称$a$为\textbf{幂零元}；否则称为非幂零元。
%\end{definition}
%
%\begin{theorem}\label{p2:prime-avoid-powers}
%若$a\in R$为非幂零元，则存在一个素理想$P$，不包含$a$的任何非负整数次幂$a^m$，$m\geq0$。
%\end{theorem}
%\begin{proof}
%令$T=\{a^m:m\geq0\}$，考虑全部与$T$不交的理想组成的集合$\mathcal S$，按包含关系排序。因为$a$非幂零，$(0)\cap T=\varnothing$，所以$\mathcal S$非空。
%
%取$\mathcal S$中的一条非空链$\{A_\alpha\}$，令$A=\bigcup_\alpha A_\alpha$。若$b,c\in A$，它们分别属于链中的两个理想；由全序性，其中一个包含另一个，所以$b,c$同属某个$A_\gamma$。于是$b-c\in A_\gamma\subseteq A$。对$r\in R$，$rb\in A$也由$b$所属理想的吸收性得到，故$A$为理想。
%若$A$包含某个$a^m$，则某个$A_\alpha$已经包含它，矛盾。所以$A\in\mathcal S$，为该链的上界。空链可取$(0)$为上界。由佐恩引理，$\mathcal S$有极大元$P$。
%
%由于$1=a^0\notin P$，$P$为真理想。反设$P$不是素理想，则存在$b,c\notin P$而$bc\in P$。令
%\[
%H=Rb+P,\qquad N=Rc+P.
%\]
%它们严格包含$P$，由$P$在$\mathcal S$中的极大性，$H,N$各与$T$相交。于是可取$a^r\in H$、$a^s\in N$，故$a^{r+s}\in HN$。
%但理想乘法展开给出
%\[
%HN=Rbc+RbP+RcP+P^2\subseteq P,
%\]
%因为$bc\in P$且其余各项也包含于$P$。于是$a^{r+s}\in P$，与$P\cap T=\varnothing$矛盾。所以$P$为所求素理想。
%\end{proof}
%
%\begin{corollary}
%每个非零交换幺环至少有一个极大理想。
%\end{corollary}
%\begin{proof}
%在上一定理的构造中取$a=1$，则$T=\{1\}$。与$T$不交的理想恰为全部真理想，因此由佐恩引理取得的极大元就是极大理想。
%\end{proof}
%
%\begin{corollary}
%$R$的全部素理想之交恰为全部幂零元组成的集合，记为
%\[
%r(R)=\bigcap_{P\text{为素理想}}P,
%\]
%称为$R$的\textbf{幂零根}。
%\end{corollary}
%\begin{proof}
%若$a^m=0$，任取素理想$P$。由$a^m\in P$反复使用素性，最终得到$a\in P$，所以$a\in r(R)$。
%若$a$非幂零，则前一定理给出不含$a$的素理想，故$a\notin r(R)$。两边相等。
%\end{proof}
%
%因此幂零元全体构成理想。也可直接验证：若$a^m=b^n=0$，则$(a-b)^{m+n-1}$的每个二项式展开项都含$a^m$或$b^n$，所以为零；对$c\in R$，$(ca)^m=c^ma^m=0$。这说明差与环元倍数仍幂零。
%
%\begin{example}
%$\Z$无非零幂零元，所以$r(\Z)=(0)$。若$n=\prod_i p_i^{e_i}>1$，由素理想对应，$\Z/(n)$中的全部素理想是$(p_i)/(n)$。因此
%\[
%r(\Z/(n))=\bigcap_i(p_i)/(n)=(p_1\cdots p_r)/(n).
%\]
%这里使用了不同素数生成的理想两两互素，从而交等于积。也可直接检查：模$n$的类$\overline a$幂零当且仅当每个$p_i$都整除$a$。
%\end{example}
%
%
%\section{商域和分式环}
%
%本节的环均为交换幺环，环之间的嵌入及以后涉及单位元的同态均保持单位元。整数环包含在有理数域中，而每个有理数都是两个整数之商。下面把这一构造推广到整环；随后再说明，存在零因子时怎样对指定的一部分元素作分母。
%
%\begin{definition}
%设$R$为整环。若域$F$包含$R$，并且$F$的每个元素都能写成
%\[
%\frac ab=ab^{-1},\qquad a,b\in R,\quad b\ne0,
%\]
%则称$F$为$R$的\textbf{商域}。
%\end{definition}
%
%\begin{theorem}\label{p2:fraction-field}
%每个整环$R$都有商域。
%\end{theorem}
%\begin{proof}
%\textbf{构造分数：}令$R^*=R\setminus\{0\}$，在$R\times R^*$中规定
%\[
%(a,b)\sim(c,d)\quad\Longleftrightarrow\quad ad=bc.
%\]
%自反性和对称性直接由等式得到。若$ad=bc$且$cf=de$，则
%\[
%adf=bcf=bde.
%\]
%由于$d\ne0$，可消去$d$，得到$af=be$。因此该关系还满足传递性，是等价关系。将$(a,b)$的等价类记作$a/b$，所有这些等价类组成的集合记作$F$。
%
%\textbf{定义运算并验证与代表元无关：}规定
%\[
%\frac ab+\frac cd=\frac{ad+bc}{bd},\qquad
%\frac ab\frac cd=\frac{ac}{bd}.
%\]
%因为$R$为整环，$bd\ne0$，所以右端仍有意义。若$ab'=a'b$、$cd'=c'd$，则
%\begin{align*}
%(ad+bc)b'd'&=(a'd'+b'c')bd,\\
%acb'd'&=a'c'bd.
%\end{align*}
%第一式可分别对两项使用$ab'=a'b$和$cd'=c'd$得到，第二式则把这两个等式相乘即可。因此两种运算均与代表元的选择无关。
%
%\textbf{验证域的运算律（补充证明）：}取$u=a/b$、$v=c/d$、$w=e/f$，有
%\begin{align*}
%(u+v)+w&=\frac{adf+bcf+bde}{bdf}=u+(v+w),\\
%(uv)w&=\frac{ace}{bdf}=u(vw),\\
%u(v+w)&=\frac{a(cf+de)}{bdf}=uv+uw.
%\end{align*}
%最后一式右端先通分为分母$bdf$。加法与乘法的交换律由$R$中的交换律得到。零元、负元和单位元分别为
%\[
%0_F=\frac01,\qquad -\frac ab=\frac{-a}{b},\qquad 1_F=\frac11.
%\]
%其中$a/b=0_F$当且仅当$a=0$，故$0_F\ne1_F$。对非零分数$a/b$，$a\ne0$，并且
%\[
%\frac ab\frac ba=\frac{ab}{ba}=1_F.
%\]
%所以$F$为域。
%
%\textbf{把$R$嵌入$F$：}映射$a\mapsto a/1$保持加法、乘法及单位元，且
%\[
%\frac a1=\frac c1\quad\Longleftrightarrow\quad a=c,
%\]
%故为单同态。把$R$与其像等同后，每个元素$a/b$恰好等于$(a/1)(b/1)^{-1}$，于是$F$是$R$的商域。
%\end{proof}
%
%\begin{theorem}\label{p2:fraction-field-extension}
%设$F$为整环$R$的商域，$\eta:R\to F'$为$R$到域$F'$的单同态。则存在唯一的单同态$\eta':F\to F'$延拓$\eta$，其表达式为
%\[
%\eta'\left(\frac ab\right)=\frac{\eta(a)}{\eta(b)}.
%\]
%\end{theorem}
%\begin{proof}
%由于$\eta$为单射，$b\ne0$蕴含$\eta(b)\ne0$，所以右端有意义。若$a/b=c/d$，则$ad=bc$，从而
%\[
%\eta(a)\eta(d)=\eta(b)\eta(c),\qquad
%\frac{\eta(a)}{\eta(b)}=\frac{\eta(c)}{\eta(d)}.
%\]
%这说明公式定义了一个映射。由分数的加法和乘法公式，
%\begin{align*}
%\eta'\left(\frac ab+\frac cd\right)
%&=\frac{\eta(a)\eta(d)+\eta(b)\eta(c)}{\eta(b)\eta(d)}
%=\eta'(a/b)+\eta'(c/d),\\
%\eta'\left(\frac ab\frac cd\right)
%&=\frac{\eta(a)\eta(c)}{\eta(b)\eta(d)}
%=\eta'(a/b)\eta'(c/d).
%\end{align*}
%又有$\eta'(1)=1$及$\eta'(a/1)=\eta(a)$，故它为所需的同态。若$\eta'(a/b)=0$，则$\eta(a)=0$，由单射性得$a=0$，所以$\eta'$也是单射。
%
%最后，任何延拓$\xi$都必须满足
%\[
%\xi(a/b)\eta(b)=\xi\bigl((a/b)b\bigr)=\eta(a),
%\]
%故$\xi(a/b)=\eta(a)/\eta(b)$。这证明唯一性。
%\end{proof}
%
%\begin{corollary}
%整环的商域在保持$R$中每个元素不动的意义下唯一。
%\end{corollary}
%\begin{proof}
%若$F_1,F_2$均为$R$的商域，对包含映射$R\hookrightarrow F_2$应用上述定理，得到单同态$F_1\to F_2$。$F_2$的每个元素都是$a/b$，而$a/b\in F_1$的像正是它，故该同态满射，是所需同构。
%\end{proof}
%
%\begin{definition}
%环$R$的非空子集$S$若对乘法封闭，就称为\textbf{乘法子集}。此处不要求$1\in S$。以后构造分式环时假定$0\notin S$。
%\end{definition}
%
%一般交换环可能有零因子。若希望$s\in S$变成可逆元，则$as=0$将强迫$a$在新环中变成零。因此先考虑
%\[
%N=\{a\in R:\text{存在 }s\in S\text{ 使 }as=0\}.
%\]
%
%\begin{lemma}
%$N$是$R$的理想。若$0\notin S$，则$N\cap S=\varnothing$；并且对$u\in S$，有
%\[
%uz\in N\quad\Longrightarrow\quad z\in N.
%\]
%\end{lemma}
%\begin{proof}
%由于$S$非空，$0\in N$。若$a_1s_1=a_2s_2=0$，其中$s_1,s_2\in S$，则
%\[
%(a_1-a_2)s_1s_2=0.
%\]
%又因$s_1s_2\in S$，$a_1-a_2\in N$。对$r\in R$，$(ra_1)s_1=r(a_1s_1)=0$，所以$ra_1\in N$，证明了理想性。
%
%若$a\in N\cap S$，存在$s\in S$使$as=0$；但乘法封闭性给出$0=as\in S$，与假设矛盾。最后，若$uz\in N$，取$t\in S$使$tuz=0$，由于$tu\in S$，即得$z\in N$。
%\end{proof}
%
%\begin{definition}
%设$S$为$R$中不含零的乘法子集。若交换幺环$R'$及保持单位元的同态$\sigma:R\to R'$满足：$\sigma(s)$对每个$s\in S$均可逆，$\ker\sigma=N$，且$R'$的每个元素均可写成$\sigma(a)\sigma(s)^{-1}$，则称$R'$为$R$关于$S$的\textbf{分式环}。
%\end{definition}
%
%\begin{theorem}\label{p2:localization-existence}
%设$R$为交换幺环，$S$为不含零的乘法子集，则$R$关于$S$的分式环存在，记作$S^{-1}R$。
%\end{theorem}
%\begin{proof}
%\textbf{等价关系：}在$R\times S$上规定
%\[
%(a,s)\sim(a',s')\quad\Longleftrightarrow\quad as'-a's\in N.
%\]
%这也等价于存在$t\in S$使$t(as'-a's)=0$。自反性和对称性由$N$为加法子群得到。若
%\[
%a_1s_2-a_2s_1\in N,\qquad a_2s_3-a_3s_2\in N,
%\]
%则
%\begin{align*}
%s_2(a_1s_3-a_3s_1)
%&=s_3(a_1s_2-a_2s_1)+s_1(a_2s_3-a_3s_2)\in N.
%\end{align*}
%应用上一引理，得$a_1s_3-a_3s_1\in N$，所以传递性成立。记等价类为$a/s$。
%
%\textbf{运算的定义：}仍规定
%\[
%\frac as+\frac bt=\frac{at+bs}{st},\qquad
%\frac as\frac bt=\frac{ab}{st}.
%\]
%分母属于$S$。为验证代表元无关，设$as'-a's\in N$、$bt'-b't\in N$。有
%\begin{align*}
%(at+bs)s't'-(a't'+b's')st
%&=tt'(as'-a's)+ss'(bt'-b't)\in N,\\
%ab s't'-a'b'st
%&=bt'(as'-a's)+a's(bt'-b't)\in N.
%\end{align*}
%因此加法、乘法均良好定义。
%
%\textbf{环的运算律（补充证明）：}若$\alpha=a/s$、$\beta=b/t$、$\gamma=c/u$，其中$s,t,u\in S$，则
%\begin{align*}
%(\alpha+\beta)+\gamma&=\frac{atu+bsu+cst}{stu}=\alpha+(\beta+\gamma),\\
%(\alpha\beta)\gamma&=\frac{abc}{stu}=\alpha(\beta\gamma),\\
%\alpha(\beta+\gamma)&=\frac{a(bu+ct)}{stu}=\alpha\beta+\alpha\gamma.
%\end{align*}
%加法与乘法均可交换。任取$t\in S$，零元、负元、单位元分别为
%\[
%0=\frac0t,\qquad -\frac as=\frac{-a}{s},\qquad 1=\frac tt.
%\]
%例如$at/(st)=a/s$，因为交叉相乘的差为零；这同时验证了上述零元和单位元的作用。不同$t$给出相同的零元和单位元。又因$t\notin N$，有$0\ne1$。
%
%\textbf{自然同态：}即使$1\notin S$，也可定义
%\[
%\sigma(a)=\frac{at}{t}\qquad(t\in S).
%\]
%换用$t'\in S$时，$att'-at't=0$，故定义与$t$无关。计算可得
%\begin{align*}
%\sigma(a)+\sigma(b)&=\frac{(a+b)t^2}{t^2}=\sigma(a+b),\\
%\sigma(a)\sigma(b)&=\frac{abt^2}{t^2}=\sigma(ab),\qquad \sigma(1)=1.
%\end{align*}
%此外，$\sigma(a)=0$当且仅当$at^2\in N$。连续应用上一引理消去$t$，此条件等价于$a\in N$。因此$\ker\sigma=N$。
%
%对$s\in S$，分数$t/(st)$是$\sigma(s)=st/t$的逆元，因为
%\[
%\frac{st}{t}\frac{t}{st}=\frac{st^2}{st^2}=1.
%\]
%并且
%\[
%\sigma(a)\sigma(s)^{-1}
%=\frac{at}{t}\frac{t}{st}
%=\frac{at^2}{st^2}=\frac as.
%\]
%所以所构造的环满足分式环定义中的全部要求。
%\end{proof}
%
%\begin{theorem}\label{p2:localization-universal}
%设$\eta:R\to R'$为交换幺环之间保持单位元的同态，并且对每个$s\in S$，$\eta(s)$可逆。则存在唯一的保持单位元的环同态
%\[
%\xi:S^{-1}R\to R',\qquad
%\xi(a/s)=\eta(a)\eta(s)^{-1},
%\]
%使$\xi\sigma=\eta$。
%\end{theorem}
%\begin{proof}
%若$a\in N$，取$t\in S$使$at=0$，得到$\eta(a)\eta(t)=0$。因$\eta(t)$可逆，$\eta(a)=0$。因此$N\subseteq\ker\eta$。
%
%若$a/s=a'/s'$，则$as'-a's\in N$，从而
%\[
%\eta(a)\eta(s')=\eta(a')\eta(s).
%\]
%乘上$\eta(s)^{-1}\eta(s')^{-1}$，即得两种代表元产生同一像，证明了良好定义。再由
%\begin{align*}
%\xi\left(\frac as+\frac bt\right)
%&=\bigl(\eta(a)\eta(t)+\eta(b)\eta(s)\bigr)
%  \bigl(\eta(s)\eta(t)\bigr)^{-1}
%=\xi(a/s)+\xi(b/t),\\
%\xi\left(\frac as\frac bt\right)
%&=\eta(a)\eta(b)\bigl(\eta(s)\eta(t)\bigr)^{-1}
%=\xi(a/s)\xi(b/t),
%\end{align*}
%可知它是环同态。还满足$\xi(t/t)=1$及$\xi(\sigma(a))=\eta(a)$。
%
%若另一同态也满足$\xi\sigma=\eta$，则由$a/s=\sigma(a)\sigma(s)^{-1}$及同态保持逆元，其在每个分数上的值都被上述公式确定，所以延拓唯一。
%\end{proof}
%
%\begin{corollary}
%分式环在与自然同态相容的意义下唯一。自然同态$\sigma:R\to S^{-1}R$为单射，当且仅当$S$中没有零因子。
%\end{corollary}
%\begin{proof}
%对两个分式环分别应用延拓定理，得到两个方向的同态；它们的复合延拓同一个自然同态，由唯一性只能是恒等映射。
%
%又因$\ker\sigma=N$，$\sigma$为单射等价于：对$s\in S$，$as=0$总蕴含$a=0$。由于$0\notin S$，这正是$S$中没有零因子的条件。特别地，对整环取$S=R\setminus\{0\}$，所得分式环就是商域。
%\end{proof}
%
%\begin{example}
%设$p$为素数，取$R=\mathbb Z$、$S=\mathbb Z\setminus(p)$。由于素数整除乘积必整除某个因子，$S$为乘法子集，并且
%\[
%S^{-1}\mathbb Z
%=\left\{\frac ab\in\mathbb Q:p\nmid b\right\}.
%\]
%该环中每个非零元素都可写成$p^t r/s$，其中$t\ge0$，$p\nmid r$、$p\nmid s$。当$t=0$时，逆元$s/r$仍在该环中；当$t>0$时，其逆元的约分后分母被$p$整除，不在该环中。所以可逆元恰为$t=0$的元素。
%
%进一步，设$H$为非零理想。在$H$的非零元素中取$p$的指数最小者$p^m r/s$。因$r/s$可逆，$p^m\in H$，故$(p^m)\subseteq H$。另一方面，$H$中每个非零元素的指数均不小于$m$，所以都属于$(p^m)$。于是
%\[
%H=(p^m).
%\]
%因而所有非零真理想为$(p),(p^2),\ldots$；其中$(p)$是唯一的极大理想，也是唯一的非零素理想。对$m\ge2$，$p\cdot p^{m-1}\in(p^m)$，但两个因子均不在$(p^m)$，故$(p^m)$不是素理想。零理想也是素理想，因为此环是整环。
%\end{example}
%
%
%\section{交换环上的多项式环}
%
%本节所有环均为交换幺环，同态保持单位元。说$R$是$R'$的子环时，本节还要求两者的单位元相同。先从扩环中的一个元素出发，说明为什么需要引入形式上的不定元。
%
%\begin{definition}
%设$R\subseteq R'$，$u\in R'$。记
%\[
%R[u]=\left\{a_0+a_1u+\cdots+a_nu^n:n\ge0,\ a_i\in R\right\},
%\qquad u^0=1.
%\]
%称它为$R$与$u$生成的子环。
%\end{definition}
%\begin{proposition}
%$R[u]$是包含$R$和$u$的最小子环。
%\end{proposition}
%\begin{proof}
%两个这样的表达式相减，仍具有相同形式。相乘时，因$u$与系数可交换，有
%\[
%\left(\sum_{i=0}^n a_i u^i\right)
%\left(\sum_{j=0}^m b_j u^j\right)
%=\sum_{k=0}^{n+m}\left(\sum_{i+j=k}a_i b_j\right)u^k.
%\]
%因此$R[u]$是子环。它包含所有常数及$u$；任何包含$R$与$u$的子环都包含$u$的各次幂及其有限线性组合，所以都包含$R[u]$。
%\end{proof}
%
%不同的系数组合可能给出$R'$中同一个元素，例如$u^2=1$时可以用$1$替换$u^2$。为把形式计算与这些特定关系分开，先构造一个系数表达式唯一的环。
%
%\begin{definition}
%以$R$中元素为系数的无穷序列
%\[
%(a_0,a_1,a_2,\ldots)
%\]
%称为一个\textbf{形式幂级数}，全体形式幂级数记为$R[[x]]$。定义逐项加法与乘法
%\begin{align*}
%(a_0,a_1,\ldots)+(b_0,b_1,\ldots)&=(a_0+b_0,a_1+b_1,\ldots),\\
%(a_0,a_1,\ldots)(b_0,b_1,\ldots)&=(c_0,c_1,\ldots),
%\qquad c_k=\sum_{i+j=k}a_i b_j.
%\end{align*}
%每个$c_k$仅涉及有限个乘积，因此定义不需要收敛性。
%\end{definition}
%
%\begin{proposition}
%$R[[x]]$在上述运算下是交换幺环。映射$a\mapsto(a,0,0,\ldots)$把$R$嵌入其中。
%\end{proposition}
%\begin{proof}
%逐项加法组成交换群，零元是全零序列，负元逐项取负。乘法可交换，因为$R$中乘法可交换。设三个序列的系数依次为$a_i,b_j,c_k$，则两种乘法结合次序所得第$n$项均为
%\[
%\sum_{i+j+k=n}a_i b_j c_k.
%\]
%这个和有限，故可重新分组，证明了结合律。类似地，第$n$项上有
%\[
%\sum_{i+j=n}a_i(b_j+c_j)
%=\sum_{i+j=n}a_i b_j+\sum_{i+j=n}a_i c_j,
%\]
%得到分配律。序列$(1,0,0,\ldots)$是单位元。常数序列的加法、乘法与$R$中相应运算一致，且不同常数给出不同序列，所以得到所述嵌入。
%\end{proof}
%
%\begin{definition}
%在$R[[x]]$中取$x=(0,1,0,\ldots)$。$R$与$x$生成的子环称为$R$上的\textbf{一元多项式环}，记为$R[x]$；$x$称为\textbf{不定元}。
%\end{definition}
%
%\begin{proposition}
%$R[x]$恰由只有有限个非零项的序列组成，且每个元素唯一地写成
%\[
%f(x)=a_0+a_1x+\cdots+a_nx^n.
%\]
%这里的唯一性指各次幂的系数唯一，允许在末尾添加零系数。
%\end{proposition}
%\begin{proof}
%按序列乘法计算，$x^i$只有第$i$个位置为$1$，其余位置均为$0$，其中第$0$个位置表示常数项。因此
%\[
%\sum_{i=0}^n a_i x^i=(a_0,a_1,\ldots,a_n,0,\ldots).
%\]
%这些序列组成子环并包含$R$和$x$；反过来每个这样的序列都是上述有限和，所以恰为$R[x]$。两个序列相等当且仅当每个位置相等，这就证明了系数的唯一性。
%\end{proof}
%
%\begin{theorem}\label{p2:poly-universal}
%设$\sigma:R\to S$为交换幺环之间保持单位元的同态，$u\in S$。则$\sigma$唯一地延拓为环同态
%\[
%\sigma_u:R[x]\to S,\qquad
%\sigma_u\left(\sum_i a_i x^i\right)=\sum_i\sigma(a_i)u^i,
%\]
%并满足$\sigma_u(x)=u$。
%\end{theorem}
%\begin{proof}
%系数表达式唯一，故公式定义了一个映射。按系数相加，得到它保持加法。对$f=\sum_i a_i x^i$、$g=\sum_j b_j x^j$，有
%\begin{align*}
%\sigma_u(fg)
%&=\sum_k\sigma\left(\sum_{i+j=k}a_i b_j\right)u^k\\
%&=\sum_{i,j}\sigma(a_i)\sigma(b_j)u^{i+j}\\
%&=\left(\sum_i\sigma(a_i)u^i\right)
%  \left(\sum_j\sigma(b_j)u^j\right)
%=\sigma_u(f)\sigma_u(g).
%\end{align*}
%所有和均为有限和，且$S$可交换，因而上述分组合法。又有$\sigma_u(1)=1$、$\sigma_u(a)=\sigma(a)$和$\sigma_u(x)=u$。
%
%若$\tau$也是这样的延拓，同态性迫使
%\[
%\tau\left(\sum_i a_i x^i\right)
%=\sum_i\tau(a_i)\tau(x)^i
%=\sum_i\sigma(a_i)u^i.
%\]
%所以$\tau=\sigma_u$，唯一性得证。
%\end{proof}
%
%\begin{definition}
%对$R\subseteq R'$和$u\in R'$，取包含映射的上述延拓，得到\textbf{代入同态}
%\[
%\varepsilon_u:R[x]\to R[u],\qquad f(x)\longmapsto f(u).
%\]
%若$f(u)=0$，称$u$为$f$的一个\textbf{根}。若存在非零$f\in R[x]$使$f(u)=0$，称$u$在$R$上\textbf{代数}；否则称$u$在$R$上\textbf{超越}。
%\end{definition}
%
%\begin{corollary}
%令$I=\ker\varepsilon_u$，则
%\[
%R[u]\cong R[x]/I,\qquad I\cap R=\{0\}.
%\]
%特别地，$u$在$R$上超越当且仅当代入同态为同构。
%\end{corollary}
%\begin{proof}
%由$R[u]$的定义，代入同态满射。环的同构定理给出商环描述。对$a\in R$，$\varepsilon_u(a)=a$，故常数的核只有零，即$I\cap R=\{0\}$。最后，超越性的定义恰好是$I=\{0\}$。
%\end{proof}
%
%\begin{proposition}
%反过来，若$I$为$R[x]$的理想且$I\cap R=\{0\}$，则$R$可嵌入$R[x]/I$。在此嵌入下，记$u=x+I$，有$R[x]/I=R[u]$。
%\end{proposition}
%\begin{proof}
%映射$a\mapsto a+I$在$R$上的核为$I\cap R$，所以为单同态。商环中每个元素可写成
%\[
%f(x)+I=(a_0+I)+(a_1+I)u+\cdots+(a_n+I)u^n,
%\]
%故整个商环由所嵌入的$R$与$u$生成。
%\end{proof}
%
%\begin{definition}
%归纳地定义\textbf{多元多项式环}
%\[
%R[x_1,\ldots,x_n]=\bigl(R[x_1,\ldots,x_{n-1}]\bigr)[x_n].
%\]
%其中$x_1^{i_1}\cdots x_n^{i_n}$称为一个\textbf{单项式}，$i_j$均为非负整数。多元多项式唯一地写成有限和
%\[
%f=\sum_{i_1,\ldots,i_n}a_{i_1\ldots i_n}x_1^{i_1}\cdots x_n^{i_n}.
%\]
%\end{definition}
%
%上述唯一性可以沿变量个数归纳：先按$x_n$的幂展开，每个系数属于$R[x_1,\ldots,x_{n-1}]$；再将这些系数按前$n-1$个变量展开。每一步的一元展开都唯一，因而最终每个单项式的系数唯一。
%
%\begin{theorem}\label{p2:multi-poly-universal}
%设$\sigma:R\to S$为交换幺环之间保持单位元的同态，$u_1,\ldots,u_n\in S$。存在唯一的环同态
%\[
%\widetilde\sigma:R[x_1,\ldots,x_n]\to S
%\]
%延拓$\sigma$且把$x_j$映为$u_j$。具体地，
%\[
%\widetilde\sigma(f)=
%\sum_{i_1,\ldots,i_n}\sigma(a_{i_1\ldots i_n})
%u_1^{i_1}\cdots u_n^{i_n}.
%\]
%\end{theorem}
%\begin{proof}
%对$n$归纳。$n=1$时即定理\ref{p2:poly-universal}。设前$n-1$个变量已经给出唯一同态
%\[
%\sigma_{n-1}:R[x_1,\ldots,x_{n-1}]\to S.
%\]
%再对系数环$R[x_1,\ldots,x_{n-1}]$应用一元情形，并指定$x_n\mapsto u_n$，便得到所需延拓。先展开$x_n$，再代入其系数，恰得到所列公式。
%
%任一满足条件的同态限制在前$n-1$个变量的多项式环上，必等于$\sigma_{n-1}$；再由一元情形的唯一性，它在加入$x_n$后也唯一。这完成归纳。
%\end{proof}
%
%\begin{definition}
%若$R\subseteq S$，记$R[u_1,\ldots,u_n]$为包含$R$及$u_1,\ldots,u_n$的最小子环，称为由这些元素生成的$R$代数。能由有限个元素这样生成的环称为\textbf{有限生成$R$代数}。
%
%若不存在非零$f\in R[x_1,\ldots,x_n]$使$f(u_1,\ldots,u_n)=0$，就称$u_1,\ldots,u_n$在$R$上\textbf{代数无关}。
%\end{definition}
%
%\begin{corollary}
%代入同态给出
%\[
%R[u_1,\ldots,u_n]\cong R[x_1,\ldots,x_n]/I,
%\quad I=\{f:f(u_1,\ldots,u_n)=0\},\quad I\cap R=\{0\}.
%\]
%这些元素代数无关当且仅当$I=\{0\}$。
%\end{corollary}
%\begin{proof}
%代入同态的像恰由所有有限多项式表达式组成，因此就是生成子环；其核为所列$I$。对常数的限制是包含映射，所以$I\cap R=\{0\}$。最后一项直接由代数无关的定义得到。
%\end{proof}
%
%\begin{proposition}
%设$N$为$R$的理想，记$N[x_1,\ldots,x_n]$为所有系数均在$N$中的多项式。则它是多项式环的理想，且
%\[
%R[x_1,\ldots,x_n]/N[x_1,\ldots,x_n]
%\cong(R/N)[y_1,\ldots,y_n].
%\]
%\end{proposition}
%\begin{proof}
%把自然同态$R\to R/N$延拓，并令$x_i\mapsto y_i$，得到
%\[
%\sum a_{i_1\ldots i_n}x_1^{i_1}\cdots x_n^{i_n}
%\longmapsto
%\sum(a_{i_1\ldots i_n}+N)y_1^{i_1}\cdots y_n^{i_n}.
%\]
%每个商环系数都可选取$R$中的代表元，故此映射满射。由多项式系数表达式的唯一性，像为零当且仅当各系数都在$N$中，其核就是$N[x_1,\ldots,x_n]$。核为理想，同构定理给出结论。
%\end{proof}
%
%生成元说明环中的元素怎样表示，而$I$说明这些表示之间有哪些关系。给出一个生成环的具体描述时，既要列出生成元，也要说明它们满足的关系；是否能用有限条关系生成整个理想，还需要关于系数环的进一步条件。
%
%
%\section{整环上的一元多项式环}
%
%\begin{definition}
%设$f(x)=a_0+\cdots+a_nx^n$且$a_n\ne0$。称$n$为$f$的\textbf{次数}，记作$\deg f$；称$a_n$为\textbf{首项系数}。首项系数为$1$的非零多项式称为\textbf{首一多项式}。约定$\deg0=-\infty$。
%\end{definition}
%
%\begin{proposition}
%对交换幺环$R$上的多项式，有
%\[
%\deg(f+g)\le\max\{\deg f,\deg g\},\qquad
%\deg(fg)\le\deg f+\deg g.
%\]
%若$\deg f\ne\deg g$，第一式取等号。若$f,g$非零且首项系数之积非零，第二式取等号；特别地，当$R$为整环时，非零多项式乘积的次数等于次数之和。
%\end{proposition}
%\begin{proof}
%两多项式相加不会出现高于两者次数的项。若次数不同，次数较高者的首项没有同次项与之抵消，所以和仍有该次数。
%
%设$f$、$g$的次数分别为$n,m$，首项系数分别为$a_n,b_m$。乘积中次数超过$n+m$的项均不存在，而$x^{n+m}$的系数恰为$a_nb_m$。因此当且仅当这个乘积非零时，乘积次数为$n+m$；否则次数下降或乘积为零。这也说明整环条件在这里起到什么作用。
%\end{proof}
%
%\begin{proposition}
%若$R$为整环，则$R[x]$也是整环，且$R[x]$中的可逆元恰为$R$中的可逆元。进而，$R[x_1,\ldots,x_n]$也是整环。
%\end{proposition}
%\begin{proof}
%两个非零多项式的首项系数均非零，其积在整环中非零，所以多项式乘积非零，证明了无零因子性。常数$1$非零，故$R[x]$是整环。
%
%若$fg=1$，则$f,g$均非零，并有
%\[
%0=\deg1=\deg f+\deg g.
%\]
%两者次数均为非负整数，只能都为零，即$f,g\in R$。因此$f$已经是$R$的可逆元。反过来，$R$中可逆元的逆仍是常数，当然也在$R[x]$中。
%
%多元情形对变量个数归纳，每次把前面的多项式环作为新的整环，应用一元结论即可。
%\end{proof}
%
%\begin{corollary}
%若$P$为交换幺环$R$的素理想，则$P[x_1,\ldots,x_n]$是$R[x_1,\ldots,x_n]$的素理想。
%\end{corollary}
%\begin{proof}
%$R/P$为整环，所以$(R/P)[y_1,\ldots,y_n]$为整环。上一节的商环同构表明，多项式环模$P[x_1,\ldots,x_n]$也是整环，故该理想为素理想。
%\end{proof}
%
%下面的带余除法不要求系数环本身为域；真正需要的是除式的首项系数可以相除。
%
%\begin{theorem}\label{p2:poly-division}
%设$R$为交换幺环，$g\in R[x]$非零且首项系数可逆。则对任意$f\in R[x]$，唯一存在$q,r\in R[x]$使
%\[
%f=qg+r,\qquad \deg r<\deg g.
%\]
%其中允许$r=0$。
%\end{theorem}
%\begin{proof}
%\textbf{存在性：}记$\deg g=m$，首项系数为$b_m$。若$f=0$或$\deg f<m$，取$q=0,r=f$即可。
%
%其余情形对$n=\deg f$归纳。设$f$的首项系数为$a_n$，取
%\[
%f_1=f-a_nb_m^{-1}x^{n-m}g.
%\]
%所减项的首项恰为$a_nx^n$，故$\deg f_1<n$。若$f_1=0$或次数小于$m$，直接令它为余式；否则由归纳假设，写成$f_1=q_1g+r$且$\deg r<m$。于是
%\[
%f=\bigl(q_1+a_nb_m^{-1}x^{n-m}\bigr)g+r,
%\]
%得到了所需分解。
%
%\textbf{唯一性：}若$f=qg+r=q'g+r'$，则
%\[
%(q-q')g=r'-r.
%\]
%若$q-q'\ne0$，它的非零首项系数乘以可逆元$b_m$仍非零。因此
%\[
%\deg\bigl((q-q')g\bigr)=\deg(q-q')+m\ge m,
%\]
%而右端次数小于$m$，矛盾。所以$q=q'$，进而$r=r'$。
%\end{proof}
%
%\begin{corollary}
%对$c\in R$及$f\in R[x]$，存在唯一的$q\in R[x]$使
%\[
%f(x)=q(x)(x-c)+f(c).
%\]
%特别地，$c$是$f$的根当且仅当$(x-c)\mid f$。
%\end{corollary}
%\begin{proof}
%用首一多项式$x-c$作除式，余式次数小于$1$，所以余式为常数$r$。代入$x=c$即得$r=f(c)$。因而余式为零恰好等价于$f(c)=0$，也恰好等价于整除关系成立。
%\end{proof}
%
%\begin{theorem}\label{p2:poly-roots}
%整环$R$上的非零$n$次多项式至多有$n$个不同的根。
%\end{theorem}
%\begin{proof}
%先把$R$嵌入其商域$F$。原多项式在$F[x]$中仍非零且次数不变，$R$中的根也都是$F$中的根，所以只须在域上证明。
%
%对次数$n$归纳。$n=0$时，非零常数没有根。若$n>0$且$f$没有根，结论成立。若有根$c$，由因式定理，
%\[
%f(x)=(x-c)g(x),\qquad\deg g=n-1.
%\]
%若$d\ne c$也是$f$的根，则
%\[
%0=f(d)=(d-c)g(d).
%\]
%由于$d-c\ne0$且$F$为域，有$g(d)=0$。除$c$之外的根均为$g$的根，由归纳假设至多有$n-1$个，加上$c$后至多有$n$个。
%\end{proof}
%
%\begin{example}
%上述条件不能任意删去。在四元数体中，若
%\[
%u=aI+bJ+cK,\qquad a,b,c\in\mathbb R,\quad a^2+b^2+c^2=1,
%\]
%则由$I^2=J^2=K^2=-1$及不同虚单位相乘反交换，$u^2=-1$。所以$x^2+1$有无穷多个四元数根；这里乘法不交换。
%
%在交换环$\mathbb Z/(8)$中，$x^2-1$有$1,3,5,7$四个根，因为奇数$2k+1$的平方减$1$为$4k(k+1)$，被$8$整除。在$\mathbb Z/(15)$中，它的根为$1,4,11,14$；可分别模$3$和模$5$选择$1$或$-1$，再用中国剩余定理得到这四个解。这两个环都有零因子，不能应用整环上的根数定理。
%\end{example}
%
%\begin{theorem}\label{p2:finite-multiplicative-subgroup}
%整环中非零元素的任一有限乘法子群都是循环群。特别地，有限域$F$的乘法群$F\setminus\{0\}$是循环群。
%\end{theorem}
%\begin{proof}
%设$G$为这样的有限子群。对任意正整数$d$，群中满足$g^d=1$的元素，都是多项式$x^d-1$的根。由根数定理，这样的元素至多有$d$个。
%
%第二章关于方程$x^d=e$的根数与循环群的判别定理\ref{p2:cyclic-root-count}说明，满足上述条件的有限群是循环群，故$G$为循环群。有限域的全部非零元素对乘法及取逆封闭，构成有限群，因而也是循环群。
%\end{proof}
%
%\begin{corollary}
%设有限域$F$有$q$个元素，则每个$a\in F$满足$a^q=a$。若$p$为素数，则模$p$有原根，原根的个数为$\varphi(p-1)$。
%\end{corollary}
%\begin{proof}
%若$a\ne0$，由乘法群的阶为$q-1$及拉格朗日定理，$a^{q-1}=1$；$a=0$时结论同样成立。
%
%所谓模$p$的原根，是指$\mathbb Z/(p)$的非零剩余类乘法群的生成元。这个群是$p-1$阶循环群，故生成元存在。循环群中生成元的个数等于欧拉函数$\varphi(p-1)$，得到计数结论。
%\end{proof}
%
%\begin{definition}
%若整环$R$的每个理想都由一个元素生成，则称$R$为\textbf{主理想整环}。例如整数环$\mathbb Z$是主理想整环。
%\end{definition}
%
%\begin{theorem}\label{p2:field-poly-pid}
%若$F$为域，则$F[x]$为主理想整环。它的每个非零理想都有唯一的首一生成元。
%\end{theorem}
%\begin{proof}
%$F[x]$为整环。零理想等于$(0)$，只需考虑非零理想$I$。在$I$中选取次数最小的非零多项式$f$。对任意$g\in I$作带余除法，
%\[
%g=qf+r,\qquad\deg r<\deg f.
%\]
%因$I$为理想，$r=g-qf\in I$。若$r\ne0$，就违反$f$次数最小的选择，所以$r=0$。于是$g\in(f)$，证明$I\subseteq(f)$；另一包含关系由$f\in I$得到。因此$I=(f)$。
%
%把$f$除以非零首项系数，可使其首一而不改变生成理想。若两个首一多项式$f,g$生成同一非零理想，则$f=ag$且$g=bf$。比较次数得$\deg f=\deg g$，故$a,b$均为非零常数。再比较首项系数得$a=b=1$，所以$f=g$。
%\end{proof}
%
%\begin{definition}
%设$F$为域，$u$是包含$F$的交换幺环中的元素，且$u$在$F$上代数。代入同态$F[x]\to F[u]$的核是非零理想，其唯一的首一生成元称为$u$在$F$上的\textbf{最小多项式}。
%\end{definition}
%
%\begin{proposition}
%设$f$为$u$的最小多项式，则$\deg f\ge1$，且对每个$g\in F[x]$，
%\[
%g(u)=0\quad\Longleftrightarrow\quad f\mid g.
%\]
%因此$f$是使$u$为根的非零多项式中次数最小的首一多项式，且
%\[
%F[u]\cong F[x]/(f).
%\]
%\end{proposition}
%\begin{proof}
%非零常数代入后不为零，所以核不含非零常数，$\deg f\ge1$。由核等于$(f)$，$g(u)=0$等价于$g\in(f)$，也就是$f\mid g$。非零倍式的次数不小于$\deg f$，所以它确有最小次数。商环同构由代入同态及同构定理得到。
%\end{proof}
%
%\begin{definition}
%域$F$上的正次数多项式$f$，若不能写成两个正次数多项式之积，就称为\textbf{不可约多项式}；否则称为\textbf{可约多项式}。等价地，$f=gh$时至少有一个因子是非零常数，即$F[x]$的可逆元。
%\end{definition}
%
%\begin{theorem}\label{p2:irreducible-maximal}
%设$F$为域，$f\in F[x]$的次数为正。以下条件等价：
%\begin{enumerate}
%\item $f$不可约；
%\item $(f)$为极大理想；
%\item $F[x]/(f)$为域；
%\item $F[x]/(f)$为整环；
%\item $(f)$为素理想。
%\end{enumerate}
%\end{theorem}
%\begin{proof}
%\textbf{由第一项推出第二项：}若$(f)\subseteq I\subseteq F[x]$，由主理想性，写$I=(g)$。因$f\in(g)$，存在$h$使$f=gh$。由不可约性，$g,h$中至少一个为可逆常数。
%
%若$g$可逆，则$I=F[x]$；若$h$可逆，则$(g)=(f)$，所以$I=(f)$。又因$\deg f>0$，$(f)$为真理想，因此它是极大理想。
%
%\textbf{第二项至第五项：}极大理想的商环为域，域是整环，而商环为整环等价于所模理想为素理想。依次应用这些已经证明的判别，得到第二项推出第三项，第三项推出第四项，第四项推出第五项。
%
%\textbf{由第五项推出第一项：}若$f$可约，写$f=gh$，其中$0<\deg g,\deg h<\deg f$。于是$gh\in(f)$。素理想性迫使$g\in(f)$或$h\in(f)$；但非零的$f$的倍式次数至少为$\deg f$，两者均不可能。矛盾，故$f$不可约。
%\end{proof}
%
%\begin{corollary}
%若$u$在域$F$上代数，最小多项式为$f$，则
%\[
%f\text{ 不可约}\quad\Longleftrightarrow\quad F[u]\text{ 为域}
%\quad\Longleftrightarrow\quad F[u]\text{ 为整环}.
%\]
%特别地，若整环$S$包含域$F$，且$S$中每个元素都在$F$上代数，则$S$本身是域。
%\end{corollary}
%\begin{proof}
%第一部分把$F[u]\cong F[x]/(f)$代入上述等价条件即可。对第二部分，任取$0\ne u\in S$，子环$F[u]$为整环，所以它已经是域。于是$u$在$F[u]\subseteq S$中有逆元。$S$中每个非零元素均可逆，故$S$为域。
%\end{proof}
%
%\begin{example}
%设$F$为有限域。由于其加法群有限，单位元的加法阶是某个正整数$p$；这就是$F$的特征。若$p=ab$且$1<a,b<p$，则
%\[
%(a\cdot1)(b\cdot1)=p\cdot1=0,
%\]
%而两个因子均非零，与域无零因子矛盾。因此$p$为素数，$F$包含由$1$生成的子域$F_p\cong\mathbb Z/(p)$。
%
%由$F$的非零乘法群循环，选其生成元$\alpha$。每个非零元素都是$\alpha$的幂，所以$F=F_p[\alpha]$。因$F$有限，$1,\alpha,\alpha^2,\ldots$中必有重复，从而$\alpha$满足非零多项式。设其最小多项式为$f$，$\deg f=n$，则
%\[
%F\cong F_p[x]/(f),\qquad f\text{ 为首一不可约多项式}.
%\]
%由带余除法，每个剩余类唯一地由次数小于$n$的多项式代表；$n$个系数各有$p$种选择，所以$|F|=p^n$。
%
%反过来，给定$F_p$上的$n$次首一不可约多项式$f$，商环$F_p[x]/(f)$是具有$p^n$个元素的域。这里尚未证明对每个$n$这样的不可约多项式都存在。
%\end{example}
%
%\begin{example}
%设$A\in M_n(F)$。把$a\in F$视为标量矩阵$aE$，则它与$A$可交换。因此即使矩阵环本身不交换，
%\[
%F[A]=\{a_0E+a_1A+\cdots+a_mA^m:a_i\in F\}
%\]
%仍是交换幺环，可以应用上述结论。
%
%矩阵空间$M_n(F)$的维数为$n^2$，所以$E,A,\ldots,A^{n^2}$线性相关，存在不全为零的$c_i$使
%\[
%c_0E+c_1A+\cdots+c_{n^2}A^{n^2}=0.
%\]
%故$A$在$F$上代数，最小多项式的次数不超过$n^2$。若最小多项式为$f$，则$F[A]\cong F[x]/(f)$；它是域当且仅当$f$不可约。
%\end{example}
%
%\begin{example}
%设$A=\begin{pmatrix}a&b\\c&d\end{pmatrix}\in M_2(\mathbb Q)$不是标量矩阵。直接计算可得
%\[
%A^2-(a+d)A+(ad-bc)E=0.
%\]
%例如左上角为$a^2+bc-(a+d)a+ad-bc=0$，右上角为$ab+bd-(a+d)b=0$；另两项同样分别为$ca+dc-(a+d)c=0$和$cb+d^2-(a+d)d+ad-bc=0$。
%
%若$A$满足一次关系$uA+vE=0$且$u\ne0$，则$A=-v/u\,E$为标量矩阵，矛盾。因此$A$的最小多项式为
%\[
%f(x)=x^2-(a+d)x+ad-bc.
%\]
%按$f$在$\mathbb Q$上的分解，有以下三种完备情形。
%
%\textbf{第一种：$f$不可约。}此时$\mathbb Q[A]$为域。每个元素唯一地写成$uE+vA$，所以它作为$\mathbb Q$上的线性空间维数为$2$。
%
%\textbf{第二种：$f=(x-\lambda_1)(x-\lambda_2)$，$\lambda_1\ne\lambda_2$。}令
%\[
%B=\frac{A-\lambda_1E}{\lambda_2-\lambda_1},\qquad
%C=-\frac{A-\lambda_2E}{\lambda_2-\lambda_1}.
%\]
%因$A$不是标量矩阵，$B,C$均非零。由$f(A)=0$，
%\[
%B+C=E,\qquad BC=CB=0.
%\]
%于是$B=B(B+C)=B^2$，同理$C^2=C$。又有$A=\lambda_2B+\lambda_1C$，所以
%\[
%\mathbb Q[A]=\mathbb QB+\mathbb QC.
%\]
%若$rB+sC=0$，分别乘$B,C$即得$rB=0,sC=0$，故$r=s=0$。两部分相乘为零，各自以$B,C$为单位元，因而
%\[
%\mathbb Q[A]=\mathbb QB\oplus\mathbb QC\cong\mathbb Q\oplus\mathbb Q.
%\]
%
%\textbf{第三种：$f=(x-\lambda)^2$。}令$B=A-\lambda E$，则$B\ne0$、$B^2=0$，并且$\mathbb Q[A]=\mathbb QE+\mathbb QB$。若$u\ne0$，则
%\[
%(uE+vB)\left(u^{-1}E-vu^{-2}B\right)=E,
%\]
%所以$uE+vB$可逆。若$u=0$，则$(vB)^2=0$，不能可逆。因此任一真理想只能包含$\mathbb QB$中的元素；若又非零，含某个$vB\ne0$，乘$v^{-1}E$得到$B$，故该理想就是$\mathbb QB$。
%
%因此它只有一个非零真理想$N=\mathbb QB$，且$N^2=0$。一般地，若理想$N$满足某个正整数$m$使$N^m=0$，便称其为\textbf{幂零理想}；这里$N$就是一个幂零理想。
%\end{example}
%
%
%\section{多项式函数}
%
%形式多项式由系数序列定义，而代入后得到的是函数。两种对象之间存在自然联系，但能否由函数值恢复多项式，需要分别讨论无限域和有限域。
%
%\begin{definition}
%设$F$为域，$F^F$表示所有从$F$到$F$的函数组成的集合。对$f,g\in F^F$，逐点定义
%\[
%(f+g)(a)=f(a)+g(a),\qquad(fg)(a)=f(a)g(a).
%\]
%常值函数$a\mapsto c$记作$f_c$。再定义函数$s\in F^F$为$s(a)=a$。
%\end{definition}
%
%\begin{proposition}
%$F^F$是交换幺环，$c\mapsto f_c$将$F$嵌入其中。子环$F[s]$中的元素称为$F$上的\textbf{多项式函数}；它们恰是形如
%\[
%a\longmapsto c_0+c_1a+\cdots+c_na^n
%\]
%的函数。
%\end{proposition}
%\begin{proof}
%函数相等当且仅当每一点的取值相等。逐点使用$F$中的加法群性质、乘法结合律及分配律，就得到函数上的相应运算律。零元与单位元分别为$f_0,f_1$，负元为$a\mapsto-f(a)$，乘法也可交换。
%
%常值函数的运算与其常数的运算一致；两个不同常数给出不同函数，故常值函数构成$F$的一个副本。最后，按生成子环的定义，$F[s]$的元素恰为有限和$\sum c_i s^i$；在$a$处取值即得所列公式。
%\end{proof}
%
%由一元多项式环的延拓定理，有满同态
%\[
%\sigma:F[x]\longrightarrow F[s],\qquad
%\sum_i c_i x^i\longmapsto\sum_i c_i s^i.
%\]
%它把形式不定元$x$映为函数$s$。记其核为$I$，则
%\[
%I=\{f\in F[x]:f(a)=0\text{ 对所有 }a\in F\text{ 成立}\},
%\qquad F[s]\cong F[x]/I.
%\]
%所以问题归结为：哪些多项式在整个域上恒等于零？
%
%\begin{theorem}\label{p2:poly-functions-one}
%若$F$为无限域，则$I=\{0\}$，从而$F[x]\cong F[s]$，每个多项式函数由唯一的形式多项式表示。
%
%若$F$为$q$元有限域，则
%\[
%I=(x^q-x),\qquad F[s]\cong F[x]/(x^q-x).
%\]
%此时每个从$F$到$F$的函数都是多项式函数，并且有唯一的次数小于$q$的多项式代表元。
%\end{theorem}
%\begin{proof}
%\textbf{无限域：}若非零$f\in I$，设$\deg f=n$，则$f$在无限多个点取零值。但根数定理说明非零$n$次多项式至多有$n$个根，矛盾。因此$I=0$，满同态$\sigma$为同构。
%
%\textbf{有限域中的核：}每个$a\in F$满足$a^q=a$，所以$x^q-x\in I$。另一方面，
%\[
%\prod_{a\in F}(x-a)=x^q-x.
%\]
%为证明此式，左边是首一$q$次多项式，并且两边在全部$q$个点取值为零。两者之差次数小于$q$却有$q$个不同的根，只能是零多项式。
%
%对$f\in I$，用$x^q-x$作带余除法，写成
%\[
%f=h(x)(x^q-x)+r(x),\qquad\deg r<q.
%\]
%代入任意$a\in F$，得到$r(a)=0$。因此$r$有$q$个根而次数小于$q$，只能为零。所以$f\in(x^q-x)$；结合反向包含关系，得$I=(x^q-x)$。
%
%\textbf{所有函数的表示：}带余除法说明每个剩余类有唯一的次数小于$q$的代表元。这些多项式共有$q^q$个，因为$q$个系数各有$q$种选择。不同代表元给出不同函数，否则差属于$I$且次数小于$q$，只能为零。
%
%从$q$元集合到自身的函数也恰有$q^q$个。因此多项式函数已包含全部函数，且每个函数的上述代表元唯一。若不限制次数，则可以加上任意$(x^q-x)$的倍式，表示便不再唯一。
%\end{proof}
%
%\begin{definition}
%对正整数$n$，令$S=F^n$，$F^S$表示从$S$到$F$的全部函数，仍按逐点运算组成交换幺环。定义第$i$个\textbf{坐标函数}
%\[
%s_i(a_1,\ldots,a_n)=a_i\qquad(1\le i\le n).
%\]
%常值函数与$s_1,\ldots,s_n$生成的子环$F[s_1,\ldots,s_n]$称为$n$元多项式函数环。
%\end{definition}
%
%多元延拓定理给出满同态
%\[
%\sigma:F[x_1,\ldots,x_n]\to F[s_1,\ldots,s_n],
%\qquad f\longmapsto f(s_1,\ldots,s_n).
%\]
%其核为在$F^n$上处处取零值的多项式组成的理想。
%
%\begin{theorem}\label{p2:poly-functions-many}
%若$F$为无限域，则$s_1,\ldots,s_n$在$F$上代数无关，且
%\[
%F[x_1,\ldots,x_n]\cong F[s_1,\ldots,s_n].
%\]
%若$F$为$q$元有限域，则
%\[
%\ker\sigma=(x_1^q-x_1,\ldots,x_n^q-x_n).
%\]
%每个$F^n\to F$的函数都是多项式函数，并且可唯一地由一个对每个变量次数均小于$q$的多项式表示。
%\end{theorem}
%\begin{proof}
%\textbf{无限域：}对变量个数$n$归纳。$n=1$已经证明。设$f$在$F^n$上处处为零，将它按最后一个变量展开为
%\[
%f=\sum_{j=0}^m f_j(x_1,\ldots,x_{n-1})x_n^j.
%\]
%固定$a_1,\ldots,a_{n-1}\in F$，所得一元多项式
%\[
%\sum_{j=0}^m f_j(a_1,\ldots,a_{n-1})x_n^j
%\]
%对每个$x_n\in F$取零值。由一元结论，它是零多项式，故每个系数$f_j(a_1,\ldots,a_{n-1})$为零。固定的点任意，因此每个$f_j$在$F^{n-1}$上处处为零。由归纳假设，$f_j=0$，所以$f=0$。这证明了$\ker\sigma=0$，也证明坐标函数代数无关。
%
%\textbf{有限域中的降次：}记
%\[
%J=(x_1^q-x_1,\ldots,x_n^q-x_n).
%\]
%由于每个$a\in F$满足$a^q=a$，各生成元在$F^n$上处处为零，所以$J\subseteq\ker\sigma$。模$J$计算时，若$e\ge q$，则
%\[
%x_i^e-x_i^{e-q+1}=x_i^{e-q}(x_i^q-x_i)\in J.
%\]
%因$q\ge2$，每次替换使指数降低$q-1$。对每个单项式的每个指数反复进行，最终都小于$q$。因此任一$f$均可写成
%\[
%f=\sum_{i=1}^n h_i(x_1,\ldots,x_n)(x_i^q-x_i)+r,
%\]
%其中$r$对每个变量的次数均小于$q$。
%
%\textbf{降次后零函数只能由零多项式表示：}仍对$n$归纳。$n=1$由根数定理成立。对$n>1$，将这样的$r$写成
%\[
%r=\sum_{j=0}^{q-1}r_j(x_1,\ldots,x_{n-1})x_n^j.
%\]
%若$r$在$F^n$上处处为零，固定前$n-1$个坐标后，所得关于$x_n$的多项式次数小于$q$，却在全部$q$个点取零值，所以各系数$r_j(a_1,\ldots,a_{n-1})$均为零。每个$r_j$对前面的变量次数仍小于$q$，由归纳假设，$r_j=0$，从而$r=0$。
%
%若$f\in\ker\sigma$，上述分解中的$r$也给出零函数，所以$r=0$，得到$f\in J$。这证明$\ker\sigma=J$。两个降次后的代表元若给出相同函数，其差也是同类的零函数，故差为零，证明了唯一性。
%
%\textbf{计数得到全部函数：}指数都小于$q$的单项式共有$q^n$个，因而这样的多项式共有$q^{q^n}$个。$F^n$有$q^n$个点，从它到$F$的函数同样共有$q^{q^n}$个。不同降次多项式已给出不同函数，数量又相同，所以全部函数均已得到。
%\end{proof}
%
%因此，在无限域上可以由全部函数值唯一确定形式多项式；在有限域上，多项式函数对应的是模上述关系理想的剩余类。讨论系数、次数或因式分解时，应始终分清所说的是形式多项式，还是它所定义的函数。
%
%
%\section{章末习题中的知识与解法}
%% !TEX root = ../../main.tex
% 第三章习题；由 chapter3.tex 按需输入。

\subsection*{环、理想与四元数（习题1--20）}

\begin{problem}
\textbf{习题1：}设环$R$中$1-ab$可逆，证明$1-ba$可逆。
\end{problem}
\begin{proof}[补充证明]
令$c=(1-ab)^{-1}$。由$c(1-ab)=(1-ab)c=1$，直接计算得到
\begin{align*}
(1-ba)(1+bca)&=1+b(c-1-abc)a=1,\\
(1+bca)(1-ba)&=1+b(c-1-cab)a=1.
\end{align*}
因此$(1-ba)^{-1}=1+b(1-ab)^{-1}a$。
\end{proof}

\begin{problem}
\textbf{习题2：}若$u$有右逆，证明以下条件等价：$u$有多于一个右逆；$u$是零因子；$u$不是单位。
\end{problem}
\begin{proof}[补充证明]
取$uv=1$。若$uw=1$且$w\ne v$，则$u(w-v)=0$，所以$u$是左零因子。若$u$是零因子，由$uv=1$可知$xu=0$蕴含$x=xuv=0$；于是存在非零$t$使$ut=0$，而$u(v+t)=1$，故右逆不唯一。

若$u$是单位，则右逆必为其唯一的逆。反之，若右逆$v$唯一，则
\[
u(vu-1)=0,
\]
故$u(v+vu-1)=1$。唯一性给出$vu=1$，从而$u$是单位。
\end{proof}

\begin{problem}
\textbf{习题3：}若$u$有多于一个右逆，证明它有无穷多个右逆。
\end{problem}
\begin{proof}[补充证明]
取$uv=1$。由于右逆不唯一，习题2说明$u$不是单位，故$e=1-vu\ne0$。又有$ue=0$和$ev=0$。对每个$k\geqslant0$，
\[
u(v+eu^k)=1.
\]
若$i<j$且$eu^i=eu^j$，右乘$v^j$，则左边为$ev^{j-i}=0$，右边为$e$，矛盾。因此这些右逆两两不同。
\end{proof}

\begin{problem}
\textbf{习题4：}证明交换环中全部幂零元构成理想，并判断非交换环中的情形。
\end{problem}
\begin{proof}[补充证明]
若$a^m=b^n=0$且$ab=ba$，则二项式展开中$(a+b)^{m+n-1}$的每一项均为零。任取$r\in R$，有$(ra)^m=r^ma^m=0$。故幂零元集对加法、取负和乘以环元封闭。

非交换情形不成立。例如在$M_2(\mathbb Z)$中，$E_{12}$和$E_{21}$均幂零，而$E_{12}+E_{21}$的平方为$1$。
\end{proof}

\begin{problem}
\textbf{习题5：}设$F$的特征为素数$p$，证明$(a+b)^p=a^p+b^p$。并将特征的讨论推广到幺环及不必有幺元的环。
\end{problem}
\begin{proof}[补充证明]
由于$F$交换，二项式公式和$p\mid\binom pk$（$0<k<p$）给出
\[
(a+b)^p=a^p+\sum_{k=1}^{p-1}\binom pk a^kb^{p-k}+b^p=a^p+b^p.
\]
幺环的特征是其加法群中$1$的阶，因而为$0$或某个正整数；若幺环是整环，正特征必须为素数。不带幺元的环以加法群的指数定义特征，即最小的使$nR=0$的正整数；若不存在则特征为$0$。交换且特征为$p$时，上述恒等式同样成立。
\end{proof}

\begin{problem}
\textbf{习题6：}设$I\lhd R$。证明$M_n(I)\lhd M_n(R)$，且$M_n(R)$的每个理想均如此得到。
\end{problem}
\begin{proof}[补充证明]
矩阵乘法按行列展开，$I$为双边理想立即给出$M_n(I)$为理想。

反过来，设$J\lhd M_n(R)$，令$I$为$J$中矩阵的全部分量所组成的集合。用矩阵单位可将$A=(a_{ij})\in J$的任一分量取出：
\[
E_{ki}AE_{j\ell}=a_{ij}E_{k\ell}\in J.
\]
于是$I$对加法及左右乘以$R$中元素封闭，且$aE_{ij}\in J$对所有$a\in I$成立。因此$J=M_n(I)$。对角元取出也说明$I$唯一。
\end{proof}

\begin{problem}
\textbf{习题7：}若$I,H,N$为$R$的理想，证明$(H+N)I=HI+NI$及$I(H+N)=IH+IN$。
\end{problem}
\begin{proof}[补充证明]
理想的积由有限个乘积之和构成。对每个$h\in H$、$n\in N$、$i\in I$，分配律给出$(h+n)i=hi+ni$，因此$(H+N)I\subseteq HI+NI$。反向包含由$H,N\subseteq H+N$得到。右侧等式同理。
\end{proof}

\begin{problem}
\textbf{习题8：}设$\eta:R\twoheadrightarrow R'$为保幺满同态。判断并证明或举反例：(i) 幂零元、幂等元的像仍是相应元素；(ii) 零因子的像仍是零因子；(iii) $R$为整环则$R'$为整环；(iv) $R'$为整环则$R$为整环；(v) 单位的像仍是单位；(vi) 像为单位的元素必为单位。
\end{problem}
\begin{proof}[补充证明]
(i) 对。若$a^m=0$，则$\eta(a)^m=0$；若$e^2=e$，则$\eta(e)^2=\eta(e)$。

(ii) 错。$\mathbb Z/6\mathbb Z\twoheadrightarrow\mathbb Z/3\mathbb Z$将零因子$4$送到非零单位$1$。

(iii) 错。$\mathbb Z\twoheadrightarrow\mathbb Z/6\mathbb Z$的源是整环，目标不是。

(iv) 错。$\mathbb Z/6\mathbb Z\twoheadrightarrow\mathbb Z/2\mathbb Z$的目标是域，源不是整环。

(v) 对。若$uv=vu=1$，则$\eta(u)\eta(v)=\eta(v)\eta(u)=1'$。

(vi) 错。$\mathbb Z\twoheadrightarrow\mathbb Z/5\mathbb Z$将非单位$2$送到单位。
\end{proof}

\begin{problem}
\textbf{习题9：}证明有理系数四元数$H_0=\{a_0+a_1I+a_2J+a_3K:a_i\in\mathbb Q\}$是$H$的子体。
\end{problem}
\begin{proof}[补充证明]
四元数的乘法表表明$H_0$对加法、乘法及取负封闭，且含$1$。若$0\ne q\in H_0$，则
\[
q^{-1}=\frac{\bar q}{q\bar q}
=\frac{a_0-a_1I-a_2J-a_3K}{a_0^2+a_1^2+a_2^2+a_3^2}\in H_0.
\]
分母是非零有理数，故$H_0$是子体。
\end{proof}

\begin{problem}
\textbf{习题10：}确定四元数体$H$的中心。
\end{problem}
\begin{proof}[补充证明]
设$q=a+bI+cJ+dK$与$I,J$都交换。比较$qI=Iq$的$J,K$系数，得$c=d=0$；再比较$qJ=Jq$，得$b=0$。故$q=a\in\mathbb R$。实数与所有四元数交换，因此中心为$\mathbb R$。
\end{proof}

\begin{problem}
\textbf{习题11：}求$H$中与$I$不交换的全部元素。
\end{problem}
\begin{proof}[补充证明]
对$q=a+bI+cJ+dK$，由$I J=K$、$J I=-K$、$I K=-J$、$K I=J$，
\[
Iq-qI=2cK-2dJ.
\]
所以所求为$\{a+bI+cJ+dK:c,d\text{不同时为零}\}$。
\end{proof}

\begin{problem}
\textbf{习题12：}设$S$是$H$的子体，且对所有$0\ne d\in H$有$dSd^{-1}\subseteq S$。证明$S=H$或$S$包含于中心。
\end{problem}
\begin{proof}[补充证明]
设$S\ne H$，取$x\in H\setminus S$及任意$s\in S$。由于$x,x+1$都非零，记
\[
a=xsx^{-1}\in S,\qquad b=(x+1)s(x+1)^{-1}\in S.
\]
由$xs=ax$及$xs+s=bx+b$得到$(a-b)x=b-s$。若$a\ne b$，则$x=(a-b)^{-1}(b-s)\in S$，矛盾。因此$a=b$，继而$s=b=a$，即$x$与$s$交换。

所以$S$外的每个元素都与$S$中的每个元素交换。对$y\in S$，$x+y\notin S$，于是$x+y$也与$s$交换；相减可知$y$与$s$交换。故$S$中的元素与$H$的每个元素交换，即$S$包含于中心。
\end{proof}

\begin{problem}
\textbf{习题13：}设$I,H,N\lhd R$，若$HN\subseteq I$且$I+H=R$，证明$N\subseteq I$。
\end{problem}
\begin{proof}[补充证明]
取$1=i+h$，其中$i\in I$、$h\in H$。任取$n\in N$，则
\[
n=(i+h)n=in+hn\in I.
\]
故$N\subseteq I$。
\end{proof}

\begin{problem}
\textbf{习题14：}设$I,H,N$为交换环$R$的理想，$I\subseteq H,N$且$H+N=R$，证明$HN\supseteq I$。
\end{problem}
\begin{proof}[补充证明]
取$1=h+n$，其中$h\in H$、$n\in N$。对每个$i\in I$，由于$I\subseteq H\cap N$，有
\[
i=ih+in\in HN.
\]
故$HN\supseteq I$。这里的包含可以是相等，例如$R=F\times F$、$I=0$、$H=F\times0$、$N=0\times F$。
\end{proof}

\begin{problem}
\textbf{习题15：}在$M_n(F)$（$n>1$）中找一个子环$R$，使$R$除恒等同构外没有其他反自同构。
\end{problem}
\begin{proof}[补充解答]
取$F$的素子域$F_0$，令$R=F_0\,1_n$。任一保幺自同构必须固定$1_n$，故固定由$1_n$生成的素子域；因$R$交换，反自同构与自同构相同。因此$R$仅有恒等反自同构。
\end{proof}

\begin{problem}
\textbf{习题16：}求$R=\left\{\begin{pmatrix}a&0\\b&c\end{pmatrix}:a,b,c\in\mathbb Z\right\}$的全部理想。
\end{problem}
\begin{proof}[补充解答]
全部理想恰为
\[
I_{m,k,n}=\left\{\begin{pmatrix}a&0\\b&c\end{pmatrix}:a\in m\mathbb Z,\ b\in k\mathbb Z,\ c\in n\mathbb Z\right\},
\]
其中$m,k,n\geqslant0$，并满足$m\mathbb Z,n\mathbb Z\subseteq k\mathbb Z$（即$k\mid m,n$，将$0$的整除意义理解为集合包含）。

事实上，用$E_{11},E_{22}$从理想中的矩阵左右相乘，可分别取出三个分量，故理想必为上述形式。再用$E_{21}$左右相乘，得到两个对角分量的取值都须属于左下角分量的取值集。反之，这两个包含关系保证上述集合对任意矩阵的左右乘法封闭。
\end{proof}

\begin{problem}
\textbf{习题17：}确定$M_2(\mathbb Z/(n))$的单位群的阶。
\end{problem}
\begin{proof}[补充解答]
由中国剩余定理，把$n=\prod p_i^{e_i}$分解，则
\[
\operatorname{GL}_2(\mathbb Z/(n))\cong\prod_i\operatorname{GL}_2(\mathbb Z/(p_i^{e_i})).
\]
模$p$约化到$\operatorname{GL}_2(\mathbb F_p)$，每个矩阵有$p^{4(e-1)}$个提升，而矩阵可逆当且仅当其模$p$像可逆。第一列有$p^2-1$种非零选择，第二列有$p^2-p$种不在第一列张成直线中的选择。故所求阶为
\[
\prod_{p^e\parallel n}p^{4(e-1)}(p^2-1)(p^2-p).
\]
当$n=1$时，该群只有一个元素。
\end{proof}

\begin{problem}
\textbf{习题18：}设交换环$R$不含零因子且仅有有限个非零元，证明$R$是有限除环。
\end{problem}
\begin{proof}[补充证明]
取$0\ne a\in R$。乘以$a$是有限集合$R$上的单射，因而满射；存在$e\in R$使$ae=a$。任取$b\in R$，由满射性写$b=ac$，于是$eb=eac=ac=b$。因此$e$是$R$的单位元。又对任意$0\ne b$，乘以$b$同样满射，存在$c$使$bc=e$。故$R$是有限域，特别是有限除环。
\end{proof}

\begin{problem}
\textbf{习题19：}若$N\lhd R$而$I\lhd N$，问$I$是否必为$R$的理想？
\end{problem}
\begin{proof}[补充解答]
未必。令$R$为$F$上所有严格上三角$3\times3$矩阵，$N=F E_{12}+F E_{13}$。直接由矩阵单位的乘法可知$N\lhd R$，且$N^2=0$，故$I=F E_{12}\lhd N$；但$E_{12}E_{23}=E_{13}\notin I$，所以$I$不是$R$的理想。
\end{proof}

\begin{problem}
\textbf{习题20：}设$R$为交换环、$a$不可逆，证明有极大理想包含$a$。
\end{problem}
\begin{proof}[补充证明]
由$a$不可逆，$(a)\ne R$。考虑所有包含$(a)$的真理想，按包含关系排序。任一全序子集的并仍是真理想：若它含$1$，则某个成员已含$1$，矛盾。由佐恩引理，有极大元$M$。它是$R$的极大理想，并含$a$。
\end{proof}

\subsection*{根、同态与素理想（习题21--32）}

\begin{problem}
\textbf{习题21：}设$J(R)$为交换环$R$的全部极大理想之交。证明$a\in J(R)$当且仅当$1-ax$对所有$x\in R$均可逆。
\end{problem}
\begin{proof}[补充证明]
若$a\in J(R)$而$1-ax$不是单位，则它属于某个极大理想$M$。因$a\in M$，便有$1=(1-ax)+ax\in M$，矛盾。

反之，若$a\notin J(R)$，取极大理想$M$使$a\notin M$。在域$R/M$中，$a+M$非零，故存在$x\in R$使$ax+M=1+M$。于是$1-ax\in M$，不是单位，矛盾。
\end{proof}

\begin{problem}
\textbf{习题22：}若$a$满足习题21的条件，称其为强拟正则元。证明全部强拟正则元构成$J(R)$。
\end{problem}
\begin{proof}[补充证明]
由习题21，对每个$a\in R$，
\[
a\text{强拟正则}\quad\Longleftrightarrow\quad a\in\bigcap_{M\text{极大}}M=J(R).
\]
故两集合相同。
\end{proof}

\begin{problem}
\textbf{习题23：}证明有限交换环的每个素理想都是极大理想。
\end{problem}
\begin{proof}[补充证明]
若$P$为素理想，则$R/P$是有限整环。由习题18，它是域，所以$P$是极大理想。
\end{proof}

\begin{problem}
\textbf{习题24：}证明第二章第2节的定理4和定理5等价。
\end{problem}
\begin{proof}[补充证明]
定理4说$\mathbb Z$的每个子群都是$m\mathbb Z$；定理5说无限循环群与$\mathbb Z$同构，并由此描述全部循环群的子群。

由定理4，若$G=\langle g\rangle$，满同态$\varphi:\mathbb Z\to G$，$k\mapsto g^k$，把$G$的任一子群$H$的原像写成$m\mathbb Z$。于是$H=\langle g^m\rangle$；当$G$无限时$\ker\varphi=0$，故$G\cong\mathbb Z$。若$G$的阶为$n$，则$\ker\varphi=n\mathbb Z\subseteq m\mathbb Z$，所以$m\mid n$，且$|H|=n/m$。反过来，$n$的每个正因子$m$都给出唯一的子群$\langle g^m\rangle$。这给出定理5。

反过来，把定理5应用于无限循环群$(\mathbb Z,+)=\langle1\rangle$，其任一子群由某个整数$m$生成，即为$m\mathbb Z$。取$m\geqslant0$并用子群的包含关系，即得定理4。
\end{proof}

\begin{problem}
\textbf{习题25：}设$G$为复数域中全部单位根构成的乘法群。给出一个满而非单的自同态。
\end{problem}
\begin{proof}[补充解答]
取$\varphi(z)=z^2$。若$z^m=1$，则$\varphi(z)^m=1$，故$\varphi$是自同态。对任意$m$次单位根$w$，取它的一个平方根$z$，则$z^{2m}=1$，所以$z\in G$且$\varphi(z)=w$。但$\varphi(1)=\varphi(-1)=1$，它并非单射。
\end{proof}

\begin{problem}
\textbf{习题26：}设$\eta:R\to R'$满足加法同态、$\eta(1)=1'$及$\eta(aba)=\eta(a)\eta(b)\eta(a)$。证明$\eta$为若尔当同态，并推出幂与极化恒等式。
\end{problem}
\begin{proof}[补充证明]
在给定恒等式中取$b=1$，即得$\eta(a^2)=\eta(a)^2$。加法性与保幺性均已给出，因此$\eta$为若尔当同态。取三重乘积中的两个外侧变量都为$a$、中间变量为$a^m$，得到
\[
\eta(a^{m+2})=\eta(a)\eta(a^m)\eta(a).
\]
以$k=0,1$为起点分别对偶数和奇数归纳，得$\eta(a^k)=\eta(a)^k$。将$a$换成$a+b$并减去平方项，得
\[
\eta(ab+ba)=\eta(a)\eta(b)+\eta(b)\eta(a).
\]
对$abc+cba$，把$a+c$代入三重乘积恒等式并减去$a$、$c$各自的项，得到
\[
\eta(abc+cba)=\eta(a)\eta(b)\eta(c)+\eta(c)\eta(b)\eta(a).
\]
\end{proof}

\begin{problem}
\textbf{习题27：}设$\eta:R\to R'$是到无零因子环的若尔当同态。证明对任意$a,b$，必有$\eta(ab)=\eta(a)\eta(b)$或$\eta(ab)=\eta(b)\eta(a)$。
\end{problem}
\begin{proof}[补充证明]
记$A=\eta(a)$、$B=\eta(b)$、$C=\eta(ab)$。在习题26的三重乘积极化式中，依次取外侧两个变量为$ab,a$，中间变量为$b$，得到
\[
\eta(ab^2a+abab)=CBA+ABC.
\]
另一方面，$\eta(ab^2a)=AB^2A$，且$\eta(abab)=\eta((ab)^2)=C^2$。故
\[
C^2+AB^2A=CBA+ABC.
\]
移项并因式分解，
\[
(C-AB)(C-BA)=0.
\]
由于$R'$无零因子，$C=AB$或$C=BA$。
\end{proof}

\begin{problem}
\textbf{习题28：}设$\eta:R\to R'$是除环之间的加法映射，$\eta(1)=1'$，且对每对$a,b$都有$\eta(ab)=\eta(a)\eta(b)$或$\eta(b)\eta(a)$。证明$\eta$为同态或反同态。
\end{problem}
\begin{proof}[补充证明]
固定$b\in R$，记
\begin{align*}
H_b&=\{a\in R:\eta(ab)=\eta(a)\eta(b)\},\\
K_b&=\{a\in R:\eta(ab)=\eta(b)\eta(a)\}.
\end{align*}
由$\eta$的加法性，$H_b,K_b$均为$R$加法群的子群。题设说明$R=H_b\cup K_b$。一个群不能是两个真子群的并，故对每个$b$，要么$H_b=R$，要么$K_b=R$。

再令$B_+=\{b:H_b=R\}$、$B_-=\{b:K_b=R\}$。它们也是$R$加法群的子群，且$R=B_+\cup B_-$。同理，$B_+=R$或$B_-=R$。前者说明$\eta$是同态，后者说明它是反同态。
\end{proof}

\begin{problem}
\textbf{习题29：}证明除环到无零因子环的若尔当同态必为同态或反同态。
\end{problem}
\begin{proof}[补充证明]
由习题27，对每一对$a,b$，$\eta(ab)$等于$\eta(a)\eta(b)$或$\eta(b)\eta(a)$。习题28的两个加法子群论证只用到$\eta$的加法性和这逐对的二择一性质，不要求目标环为除环。因此它也适用于此处，并给出$\eta$为同态或反同态。
\end{proof}

\begin{problem}
\textbf{习题30：}设$R=\mathbb Z/(p^n)$，其中$p$为素数、$n>1$。确定零因子、素理想及相应商域。
\end{problem}
\begin{proof}[补充解答]
整数$a$在$R$中可逆当且仅当$p\nmid a$。若$p\mid a$，则$a$乘以某个非零的$p$幂即为$0$，故非单位恰为零因子。理想由$p^n$的因子对应，唯一的素理想为$P=(p)/(p^n)$。约化模$p$给出
\[
R/P\cong\mathbb Z/(p),
\]
它是域。
\end{proof}

\begin{problem}
\textbf{习题31：}设$R=\mathbb Z/(p^n)$，确定$R[x]$的单位和零因子。
\end{problem}
\begin{proof}[补充解答]
写$f=a_0+a_1x+\cdots+a_mx^m$。由习题43和45，$f$可逆当且仅当$a_0$不被$p$整除、$a_i$（$i>0$）均被$p$整除。

若非零$f$的所有系数都被$p$整除，则$p^{n-1}f=0$，故$f$是零因子。反之，设$a_j$是$f$的最低次单位系数。任取非零$g=\sum b_kx^k$，令$s$为$b_k$的最小$p$进赋值，并取具有此赋值的最低次系数$b_k$。在$fg$的$x^{j+k}$项中，$a_jb_k$的赋值恰为$s$；其余乘积或使用次数小于$j$的$f$系数，或使用次数小于$k$的$g$系数，赋值都大于$s$。该项因而非零，$fg\ne0$。所以非零零因子恰为全部系数均被$p$整除的多项式。
\end{proof}

\begin{problem}
\textbf{习题32：}设$R=R_1\oplus\cdots\oplus R_r$，$N\lhd R$。证明$N=\bigoplus_i(N\cap R_i)$，并求$\mathbb Z/(n)$的最大幂零理想和单位群。
\end{problem}
\begin{proof}[补充证明]
设$e_i$是第$i$个直和因子的单位元。任取$x\in N$，
\[
x=\sum_i e_ix,
\]
其中$e_ix\in N\cap R_i$，且分解唯一。因此得到所述直和。

若$n=\prod_i p_i^{a_i}$，中国剩余定理给出$\mathbb Z/(n)\cong\prod_i\mathbb Z/(p_i^{a_i})$。每一因子的幂零元恰为$p_i$的倍数，故最大幂零理想为$(\prod_i p_i)/(n)$；单位群为
\[
(\mathbb Z/(n))^\times\cong\prod_i(\mathbb Z/(p_i^{a_i}))^\times,
\qquad |(\mathbb Z/(n))^\times|=n\prod_{p\mid n}\left(1-\frac1p\right).
\]
\end{proof}

\subsection*{多项式与有限域（习题33--51）}

\begin{problem}
\textbf{习题33：}证明$x^3-x$在$\mathbb Z/(6)$中有六个根。
\end{problem}
\begin{proof}[补充证明]
对任意整数$a$，三个连续整数$a-1,a,a+1$的乘积同时被$2$和$3$整除。因此$6\mid a^3-a$。$\mathbb Z/(6)$的全部六个元素都是根。
\end{proof}

\begin{problem}
\textbf{习题34：}证明$x^2+1$在四元数体$H$中有无穷多个根。
\end{problem}
\begin{proof}[补充证明]
若$q=bI+cJ+dK$且$b^2+c^2+d^2=1$，四元数乘法关系使交叉项成对抵消，故
\[
q^2=-(b^2+c^2+d^2)=-1.
\]
单位球面上的三元组有无穷多个，因此根有无穷多个。
\end{proof}

\begin{problem}
\textbf{习题35：}设$u=\sqrt2+\sqrt3$，求其在$\mathbb Q$上的极小多项式。
\end{problem}
\begin{proof}[补充解答]
由$u^2=5+2\sqrt6$，得
\[
(u^2-5)^2=24,
\qquad u^4-10u^2+1=0.
\]
又$u^{-1}=\sqrt3-\sqrt2$，故$\sqrt3=(u+u^{-1})/2$、$\sqrt2=(u-u^{-1})/2$均属于$\mathbb Q(u)$。于是$\mathbb Q(u)=\mathbb Q(\sqrt2,\sqrt3)$，其次数为$4$。故极小多项式为$x^4-10x^2+1$。
\end{proof}

\begin{problem}
\textbf{习题36：}设$u^3-u+1=0$。将$(5u^2+3u-1)(2u^2-2u+6)$及$(3u^2-u+2)^{-1}$写成次数不超过$2$的多项式。
\end{problem}
\begin{proof}[补充解答]
用$u^3=u-1$、$u^4=u^2-u$逐项约去高次幂，得到
\[
(5u^2+3u-1)(2u^2-2u+6)=32u^2+6u-2.
\]
再直接相乘并约去高次幂，可验证
\[
(3u^2-u+2)(-2u^2-u+3)=7.
\]
因此$(3u^2-u+2)^{-1}=(-2u^2-u+3)/7$。
\end{proof}

\begin{problem}
\textbf{习题37：}证明$\mathbb F_2[x]/(x^3+x^2+1)$是八元域。
\end{problem}
\begin{proof}[补充证明]
三次多项式在域上不可约当且仅当它没有根。代入$0,1$分别得到$1,1$，所以$x^3+x^2+1$不可约。商环是域，每个剩余类唯一写作$a+bx+cx^2$，其中$a,b,c\in\mathbb F_2$，故有$2^3=8$个元素。
\end{proof}

\begin{problem}
\textbf{习题38：}列出$\mathbb F_5[x]$中全部首项系数为$1$的二次不可约多项式。
\end{problem}
\begin{proof}[补充解答]
$x^2+bx+c$不可约当且仅当判别式$b^2-4c$不是$\mathbb F_5$中的平方，即等于$2$或$3$。逐一代入$b=0,1,2,3,4$，得到
\begin{align*}
&x^2+2,\ x^2+3,\ x^2+x+1,\ x^2+x+2,\ x^2+2x+3,\\
&x^2+2x+4,\ x^2+3x+3,\ x^2+3x+4,\ x^2+4x+1,\ x^2+4x+2.
\end{align*}
\end{proof}

\begin{problem}
\textbf{习题39：}构造一个有$25$个元素的域。
\end{problem}
\begin{proof}[补充解答]
习题38表明$x^2+2$在$\mathbb F_5[x]$中不可约。故
\[
F=\mathbb F_5[x]/(x^2+2)
\]
是域，其元素唯一写作$a+b\bar x$（$a,b\in\mathbb F_5$），共有$5^2=25$个。
\end{proof}

\begin{problem}
\textbf{习题40：}证明$(3,x^3+2x^2+2x-1)$是$\mathbb Z[x]$的极大理想。
\end{problem}
\begin{proof}[补充证明]
模$3$约化后，多项式成为$h=x^3-x^2-x-1\in\mathbb F_3[x]$。代入$0,1,2$分别得$2,1,1$，所以三次多项式$h$无根，因而不可约。于是
\[
\mathbb Z[x]/(3,x^3+2x^2+2x-1)\cong\mathbb F_3[x]/(h)
\]
是域，故该理想极大。
\end{proof}

\begin{problem}
\textbf{习题41：}令$I=(3,x^3+2x^2+2x-1)$。证明$\mathbb Z[x]/I\cong\mathbb F_3[x]/(x^3-x^2-x-1)$，并推出$I$极大。
\end{problem}
\begin{proof}[补充证明]
系数模$3$的满同态$\rho:\mathbb Z[x]\to\mathbb F_3[x]$以$(3)$为核，并把$x^3+2x^2+2x-1$送到$x^3-x^2-x-1$。同态定理给出所述同构。由习题40中该三次式的不可约性，右侧是域，故$I$极大。
\end{proof}

\begin{problem}
\textbf{习题42：}设$F$为域且$\deg f>0$。证明$F[x]/(f)$无零因子当且仅当$f$不可约。
\end{problem}
\begin{proof}[补充证明]
若$f$不可约，则$(f)$为极大理想，商环为域。反之，若$f=gh$且$0<\deg g,\deg h<\deg f$，则$\bar g,\bar h$在商环中均非零，但$\bar g\bar h=0$，矛盾。
\end{proof}

\begin{problem}
\textbf{习题43：}设$R$为交换环。若$f\in R[x]$幂零，证明$f$的每个系数均幂零。
\end{problem}
\begin{proof}[补充证明]
写$f=a_0+a_1x+\cdots+a_mx^m$，设$f^N=0$。最高项系数为$a_m^N$，故$a_m$幂零。模去幂零理想$(a_m)$后，$f$的次数下降，且仍幂零。对次数归纳，所有$a_i$在$R/(a_m)$中幂零；于是某个$a_i^k\in(a_m)$，而$(a_m)$为幂零理想，故$a_i$也幂零。
\end{proof}

\begin{problem}
\textbf{习题44：}设$R$为交换环。证明$f\in R[x]$为零因子当且仅当存在$0\ne a\in R$使$af=0$。
\end{problem}
\begin{proof}[补充证明]
充分性直接成立。若$fg=0$且$g\ne0$，在所有被$f$消去的非零多项式中选次数最低者$g$。设$f=a_0+\cdots+a_mx^m$。因$fg=0$，最高次项系数给出$a_m$乘以$g$的最高次系数为零，故$\deg(a_mg)<\deg g$。同时$f(a_mg)=a_m(fg)=0$，由$g$的最小性得$a_mg=0$。于是$(f-a_mx^m)g=0$。依次降次，得到$a_ig=0$对每个$i$成立。取$g$的任一非零系数$a$，便有$af=0$。
\end{proof}

\begin{problem}
\textbf{习题45：}设$R$为交换环，$f=a_0x^n+a_1x^{n-1}+\cdots+a_n$。若$a_n$是单位而$a_0,\ldots,a_{n-1}$均幂零，证明$f$是$R[x]$的单位。
\end{problem}
\begin{proof}[补充证明]
令$h=a_n^{-1}(f-a_n)$。它只有有限个非零系数，且每个系数都幂零；在足够高的幂$h^N$的任一项中，必有某个系数出现到其幂零次数，故$h^N=0$。于是
\[
f^{-1}=a_n^{-1}(1+h)^{-1}
=a_n^{-1}\sum_{j=0}^{N-1}(-h)^j.
\]
\end{proof}

\begin{problem}
\textbf{习题46：}证明习题43的逆命题。
\end{problem}
\begin{proof}[补充证明]
若$f=\sum_{i=0}^m a_ix^i$且$a_i^{n_i}=0$，取$N=1+\sum_{i=0}^m(n_i-1)$。$f^N$的每一项都含$N$个系数因子；由抽屉原理，其中某个$a_i$至少出现$n_i$次，所以该项为零。故$f^N=0$。
\end{proof}

\begin{problem}
\textbf{习题47：}设$R'=R[x_1,\ldots,x_n]$，$r(R)$为$R$的幂零根。证明$r(R')=r(R)[x_1,\ldots,x_n]$。
\end{problem}
\begin{proof}[补充证明]
对变量个数归纳。习题43和46说明一元多项式幂零当且仅当各系数幂零，即$r(R[x])=r(R)[x]$。将$R[x_1,\ldots,x_{n-1}]$代替$R$，重复这一等式，便得结论。
\end{proof}

\begin{problem}
\textbf{习题48：}证明有限环中右逆也是左逆，并判断无限环是否仍成立。
\end{problem}
\begin{proof}[补充证明]
若$uv=1$且$R$有限，映射$x\mapsto ux$有右逆映射$x\mapsto vx$，故是满射，因有限而单射。由$u(vu-1)=0$得$vu=1$。无限环不成立：在可数维向量空间的自同态环中，左移算子$u$与右移算子$v$满足$uv=1$而$vu\ne1$。
\end{proof}

\begin{problem}
\textbf{习题49：}证明交换环$R$的多项式环$R[x]$中，强拟正则元恰为幂零元。
\end{problem}
\begin{proof}[补充证明]
若$f$强拟正则，则$1-xf$是$R[x]$的单位。多项式单位的正次数系数均幂零，故$f$的各系数均幂零。由习题46，$f$幂零。

反之，若$f^N=0$，则对任意$g\in R[x]$，$(fg)^N=0$，且
\[
(1-fg)^{-1}=\sum_{j=0}^{N-1}(fg)^j.
\]
由习题21的刻画，$f$强拟正则。
\end{proof}

\begin{problem}
\textbf{习题50：}设交换环$R$中每个$a$都有某个$n>1$使$a^n=a$。证明每个素理想均极大。
\end{problem}
\begin{proof}[补充证明]
若$P$为素理想，$R/P$是整环，且其中每个元素仍满足某个$\bar a^n=\bar a$。对$\bar a\ne0$，消去$\bar a$得$\bar a^{n-1}=1$，故$\bar a$可逆。于是$R/P$为域，$P$极大。
\end{proof}

\begin{problem}
\textbf{习题51：}设$R$为布尔环，即每个$a$满足$a^2=a$。证明$2a=0$、每个素理想极大且商域为$\mathbb F_2$、每个有限生成理想均为主理想。
\end{problem}
\begin{proof}[补充证明]
由$(a+a)^2=a+a$及$a^2=a$，展开并约去$2a$，得$2a=0$。又$(a+b)^2=a+b$给出$ab=ba$，故$R$交换。若$P$素，则$R/P$是布尔整环；每个非零元由$a^2=a$得$a=1$，所以$R/P\cong\mathbb F_2$，$P$极大。

对$a,b\in R$，记$c=a+b+ab$。因$ac=a$、$bc=b$，有$(a,b)\subseteq(c)$；又$c\in(a,b)$，故$(a,b)=(c)$。对生成元个数归纳，任一有限生成理想均为主理想。
\end{proof}

\subsection*{幂零根、有限域与形式幂级数（习题52--68）}

\begin{problem}
\textbf{习题52：}设$r(R)$为交换环$R$的幂零根。证明以下三项等价：$R$只有一个素理想；$R$的每个元素不是单位就是幂零元；$R/r(R)$是域。
\end{problem}
\begin{proof}[补充证明]
幂零根是全部素理想之交。若素理想只有$P$，则$r(R)=P$；任一极大理想均为素理想，故$P$是唯一极大理想。$P$外元素均为单位，$P$内元素均幂零。

若每个非单位均幂零，则$r(R)$恰为非单位集，商环中每个非零元都可逆，故$R/r(R)$是域。若商环是域，则$r(R)$极大；任一素理想都含它，只能等于它，故素理想唯一。
\end{proof}

\begin{problem}
\textbf{习题53：}证明只有有限多个理想的整环是域。
\end{problem}
\begin{proof}[补充证明]
对$0\ne a\in R$，链$(a)\supseteq(a^2)\supseteq\cdots$必有相邻两项相等。于是$a^m=a^{m+1}b$，在整环中消去$a^m$得$1=ab$。所以每个非零元可逆。
\end{proof}

\begin{problem}
\textbf{习题54：}设整环$R$的任意非空理想集合按包含关系都有极小元。证明$R$是域。
\end{problem}
\begin{proof}[补充证明]
对$0\ne a\in R$，在$\{(a^m):m\geqslant1\}$中取极小元$(a^m)$。由$(a^{m+1})\subseteq(a^m)$及极小性，二者相等。按习题53消去，$a$可逆。
\end{proof}

\begin{problem}
\textbf{习题55：}设$a,b,ab-1$均为单位。证明$a-b^{-1}$和$(a-b^{-1})^{-1}-a^{-1}$均为单位，并验证华罗庚恒等式。
\end{problem}
\begin{proof}[补充证明]
令$c=a-b^{-1}=(ab-1)b^{-1}$，故$c$可逆。又
\[
c^{-1}-a^{-1}=c^{-1}(a-c)a^{-1}=c^{-1}b^{-1}a^{-1}
\]
是单位。取逆并代入$c$，得到
\[
\bigl((a-b^{-1})^{-1}-a^{-1}\bigr)^{-1}=abc=aba-a.
\]
\end{proof}

\begin{problem}
\textbf{习题56：}设$D$为除环，$C$为其中心，$R$为子体且对每个$0\ne d\in D$有$dRd^{-1}\subseteq R$。证明$R=D$或$R\subseteq C$。
\end{problem}
\begin{proof}[补充证明]
若$R\ne D$，取$x\in D\setminus R$及$a\in R$。记$b=xax^{-1}$、$c=(x+1)a(x+1)^{-1}$，二者均在$R$。比较$xa=bx$与$xa+a=cx+c$，得$(b-c)x=c-a$。若$b\ne c$，则$x=(b-c)^{-1}(c-a)\in R$，矛盾。故$b=c=a$，即$x$与$a$交换。

任取$y\in R$，$x+y\notin R$，故$x+y$也与$a$交换，从而$y$与$a$交换。于是$a$与$D$的每个元素交换，即$a\in C$。故$R\subseteq C$。
\end{proof}

\begin{problem}
\textbf{习题57：}设$R$为实系数多项式函数环。证明每个$t\in\mathbb R$的代入映射$\eta_t(f)=f(t)$为到$\mathbb R$的同态，且每个这样的同态均为某个$\eta_t$。
\end{problem}
\begin{proof}[补充证明]
代入保持$1$、加法及乘法。设$\eta:R\to\mathbb R$为保幺同态。对实常数$c>0$，写$c=d^2$，则$\eta(c)=\eta(d)^2\geqslant0$；故$\eta$在常数域上保序。它固定有理数，由有理数的稠密性，$\eta(c)=c$对每个实数$c$成立。令$t=\eta(x)$，则
\[
\eta\left(\sum_i a_ix^i\right)=\sum_i a_it^i.
\]
\end{proof}

\begin{problem}
\textbf{习题58：}设$F$为无限域。若$g\ne0$，且$f$在所有满足$g(c_1,\ldots,c_n)\ne0$的点上为零，证明$f=0$。
\end{problem}
\begin{proof}[补充证明]
在$g$非零的点上$fg=0$，在$g$为零的点上亦然。故$fg$在$F^n$上恒为零。无限域上的多项式若处处为零，则为零多项式；由多项式环无零因子且$g\ne0$，得$f=0$。
\end{proof}

\begin{problem}
\textbf{习题59：}设$F$有$q$个元素。证明$g_0=\prod_{i=1}^n(1-x_i^{q-1})$在原点取值$1$，在其余点取值$0$。
\end{problem}
\begin{proof}[补充证明]
在原点，各因子均为$1$。若$a=(a_1,\ldots,a_n)\ne0$，某个$a_i\ne0$；由有限域乘法群的阶为$q-1$，$a_i^{q-1}=1$，相应因子为$0$。
\end{proof}

\begin{problem}
\textbf{习题60：}若$f$在原点为零而在其余点均非零，证明$g=1-f^{q-1}$与习题59的$g_0$取值相同，并推出$\deg f\geqslant n$。
\end{problem}
\begin{proof}[补充证明]
由$z^{q-1}=1$对每个$z\in F^\times$成立，$g$在原点为$1$，在其余点为$0$。将多项式逐变量模$x_i^q-x_i$约化，不增加总次数，可得到各变量次数小于$q$的唯一函数代表。$g_0$就是此代表，且其单项式$x_1^{q-1}\cdots x_n^{q-1}$系数为$(-1)^n\ne0$。因此$\deg g\geqslant n(q-1)$；又$\deg g\leqslant(q-1)\deg f$，故$\deg f\geqslant n$。
\end{proof}

\begin{problem}
\textbf{习题61：}设$F$是$q$元域，$\deg f<n$且$f(0)=0$。证明$f$还有一个非零零点。
\end{problem}
\begin{proof}[补充证明]
记$N$为$f$在$F^n$中的零点数。在$F$中计算
\[
N=\sum_{a\in F^n}\bigl(1-f(a)^{q-1}\bigr).
\]
对$k\geqslant0$，$\sum_{t\in F}t^k$只有在$k$为$q-1$的正倍数时才可能非零；当$k=0$时该和为$q=0$。故一个单项式对全体点求和后可能非零，要求每个变量的指数至少为$q-1$，总次数至少$n(q-1)$。但$\deg f^{q-1}<n(q-1)$，故$N=0$在$F$中，即整数$N$被特征$p$整除。原点已是零点，故$N\geqslant p\geqslant2$。
\end{proof}

\begin{problem}
\textbf{习题62：}设$R$为交换环、$f=\sum_{i\geqslant0}a_ix^i\in R[[x]]$。证明：(i) $f$可逆当且仅当$a_0$可逆；(ii) $R[[x]]$为整环当且仅当$R$为整环；(iii) 若$R$为整环且$F$为其商域，则$R[[x]]$商域的非零元均可写作$x^\nu(a_0+a_1x+\cdots)$，其中$\nu\in\mathbb Z$、$a_i\in F$、$a_0\ne0$。
\end{problem}
\begin{proof}[补充证明]
设$f=\sum_{i\geqslant0}a_ix^i$。若$fg=1$，常数项给出$a_0b_0=1$。反之，若$a_0$可逆，递归取
\[
b_0=a_0^{-1},\qquad b_m=-a_0^{-1}\sum_{i=1}^m a_i b_{m-i}\quad(m\geqslant1),
\]
则$g=\sum b_ix^i$满足$fg=1$。

若$R$为整环，非零幂级数的最低非零项相乘仍非零，故$R[[x]]$为整环；反向由常数嵌入可知。设$F$为$R$的商域，$R[[x]]$的商域自然嵌入$F((x))$。因此其中每个非零元均可写作
\[
x^\nu(a_0+a_1x+a_2x^2+\cdots),\qquad \nu\in\mathbb Z,\ a_i\in F,\ a_0\ne0.
\]
当$R=F$本身是域时，$F[[x]]$的商域恰为$F((x))$。
\end{proof}

\begin{problem}
\textbf{习题63：}若交换环$R$中每个非单位$a$都使$1+a$为单位，证明全部非单位构成唯一的极大理想，并判断这种环是否存在。
\end{problem}
\begin{proof}[补充证明]
记非单位集为$M$。若$a\in M$、$r\in R$而$ra$为单位，则$a$也为单位，矛盾；故$ra\in M$。若$a,b\in M$而$a+b=u$为单位，则$a/u,b/u$仍非单位。由题设，$1-a/u=b/u$应为单位，矛盾。因此$M$对加法及取负封闭，是理想。它包含所有真理想，因而是唯一极大理想。例如$\mathbb Z/(p^n)$满足该性质。
\end{proof}

\begin{problem}
\textbf{习题64：}设交换整环$R$中每个非零非单位$a$都使$1+a$仍非单位。证明全部单位加上$0$构成子域$F$，且$R\setminus F$中的元素都超越于$F$。
\end{problem}
\begin{proof}[补充证明]
记$F=R^\times\cup\{0\}$。若单位$u,v$之和非零且非单位，令$a=-(u+v)u^{-1}$，则$a$非零且非单位；但$1+a=-vu^{-1}$为单位，矛盾。所以$F$对加法封闭；其余域运算显然封闭。

若$z\in R\setminus F$在$F$上代数，取常数项非零的非零关系$c_0+c_1z+\cdots+c_mz^m=0$；若原关系常数项为零，可在整环中消去适当的$z$幂。于是
\[
z(c_1+\cdots+c_mz^{m-1})=-c_0\in F^\times,
\]
说明$z$可逆，矛盾。
\end{proof}

\begin{problem}
\textbf{习题65：}设$F$为无限域。证明不存在非常数多项式$f\in F[x]$满足$f(x)f(y)=f(x+y)$。
\end{problem}
\begin{proof}[补充证明]
设$\deg f=d>0$、首项系数为$c\ne0$。等式左边含非零项$c^2x^dy^d$，总次数为$2d$；右边总次数为$d$，矛盾。
\end{proof}

\begin{problem}
\textbf{习题66：}设$F$为域，证明$M_n(F)$的每个左理想都是主理想。
\end{problem}
\begin{proof}[补充证明]
设$L$为左理想，把$L$中各矩阵的行向量组成的子空间记为$W\subseteq F^n$。左乘矩阵单位可把任一行放入任意指定的行，其余行置零；故$L$恰是所有行向量均属于$W$的矩阵。取矩阵$A$，使其非零行构成$W$的一组基，则$L=M_n(F)A$。
\end{proof}

\begin{problem}
\textbf{习题67：}设$G$为阿贝尔群，证明$\operatorname{End}(G)$以逐点加法和映射复合构成含幺环。
\end{problem}
\begin{proof}[补充证明]
若$\eta,\xi$为自同态，则
\[
(\eta+\xi)(a+b)=\eta(a)+\eta(b)+\xi(a)+\xi(b)
=(\eta+\xi)(a)+(\eta+\xi)(b),
\]
重排使用$G$交换。零映射、逐点负映射给出加法零元与逆元；加法结合且交换。复合结合，恒等映射为乘法单位。自同态的加法性给出复合对逐点加法的左右分配律。
\end{proof}

\begin{problem}
\textbf{习题68：}分别求$\mathbb Z^+$、无限循环群、$\mathbb Q^+$及$n$个$p$阶循环群的直和的自同态环。
\end{problem}
\begin{proof}[补充解答]
$\mathbb Z^+$及任一无限循环群的自同态由生成元的像唯一决定；该像是生成元的整数倍，故自同态环均同构于$\mathbb Z$。

$\mathbb Q^+$的自同态$\eta$由$\eta(1)=r\in\mathbb Q$唯一决定，因为$n\eta(m/n)=m\eta(1)$。故$\eta(q)=rq$，自同态环同构于$\mathbb Q$。

$n$个$p$阶循环群的直和是$\mathbb F_p$上的$n$维向量空间，自同态恰为线性变换，所以其自同态环同构于$M_n(\mathbb F_p)$。
\end{proof}


%